% study-pack shared preamble (pdfLaTeX). Copy next to your .tex; don't edit it per document.
% Usage, after \documentclass and geometry:
%   \def\studylang{english}    % output language: english, russian, ngerman, french, spanish, ...
%   \def\sourcelang{english}   % language of the course materials (for \term{...}); optional
%   % study-pack shared preamble (pdfLaTeX). Copy next to your .tex; don't edit it per document.
% Usage, after \documentclass and geometry:
%   \def\studylang{english}    % output language: english, russian, ngerman, french, spanish, ...
%   \def\sourcelang{english}   % language of the course materials (for \term{...}); optional
%   % study-pack shared preamble (pdfLaTeX). Copy next to your .tex; don't edit it per document.
% Usage, after \documentclass and geometry:
%   \def\studylang{english}    % output language: english, russian, ngerman, french, spanish, ...
%   \def\sourcelang{english}   % language of the course materials (for \term{...}); optional
%   % study-pack shared preamble (pdfLaTeX). Copy next to your .tex; don't edit it per document.
% Usage, after \documentclass and geometry:
%   \def\studylang{english}    % output language: english, russian, ngerman, french, spanish, ...
%   \def\sourcelang{english}   % language of the course materials (for \term{...}); optional
%   \input{preamble}
% Box titles are built in for english, russian, ngerman/german, french, spanish. For another
% language, after \input{preamble}, \renewcommand the \ttl... macros below.
%
% Boxes: defn idea trap method task exmpl keyres srcerr.
%   \begin{trap}[: suffix] ... \end{trap}             title "Trap: suffix"
%   \begin{exmpl}[ 3][breakable] ... \end{exmpl}      long box allowed to split across pages
% Boxes do NOT split across pages by default; pass [breakable] as the 2nd option only for a box
% longer than about half a page.
\providecommand\studylang{english}
\providecommand\sourcelang{english}
\usepackage{etoolbox}
\usepackage[T1,T2A]{fontenc}
\usepackage[utf8]{inputenc}
% Spanish babel's active shorthands (and decimal-dot handling) break TikZ, \% spacing and icomma.
\ifboolexpr{test {\ifdefstring{\studylang}{spanish}} or test {\ifdefstring{\sourcelang}{spanish}}}
  {\PassOptionsToPackage{es-noshorthands,es-nodecimaldot}{babel}}{}
\ifdefstring{\studylang}{\sourcelang}
  {\usepackage[\studylang]{babel}}
  {\usepackage[\sourcelang,\studylang]{babel}}
\usepackage{paratype}
\usepackage{amsmath,amssymb,mathtools}
% Decimal comma for non-English output. In math, type a space after any comma that separates
% items: (1, 2), (x, y), a_{1, 4}. Otherwise (1,2) prints as the decimal 1,2.
\ifdefstring{\studylang}{english}{}{\usepackage{icomma}}
\usepackage{xcolor}
\usepackage[most]{tcolorbox}
\usepackage{enumitem}
\usepackage{booktabs,array,tabularx,multirow}
\usepackage{pgfplots}
\pgfplotsset{compat=1.18}
\usetikzlibrary{arrows.meta,positioning,calc,shapes.geometric,fit,backgrounds,babel}% babel: French/Turkish active : ; ! stay safe in TikZ
\usepackage{titlesec}
\usepackage{fancyhdr}
\usepackage{multicol}
\usepackage{needspace}
\usepackage[hidelinks]{hyperref}

% --- layout defaults that prevent the usual fix-up rounds ---
\setlength{\parindent}{0pt}
\setlength{\parskip}{3pt plus 1pt}
\setlength{\emergencystretch}{3em}   % no overfull text lines
\allowdisplaybreaks                   % long align blocks may break across pages
\clubpenalty=10000 \widowpenalty=10000 \displaywidowpenalty=10000
\raggedbottom                         % no stretched pages before an unbreakable box
\setlist{nosep,leftmargin=1.4em}

\definecolor{accent}{RGB}{22,70,135}
\definecolor{idea}{RGB}{0,120,110}
\definecolor{trap}{RGB}{190,80,0}
\definecolor{ex}{RGB}{60,110,30}
\definecolor{grayline}{RGB}{150,150,150}
\definecolor{c1}{RGB}{22,70,135}
\definecolor{c2}{RGB}{200,90,20}
\definecolor{c3}{RGB}{40,140,60}
\definecolor{c4}{RGB}{170,40,90}

% Source-language term inside output-language text: \term{deadweight loss}
\newcommand{\term}[1]{\textit{\foreignlanguage{\sourcelang}{#1}}}

% --- box titles ---
\newcommand\ttlDefinition{Definition}\newcommand\ttlIntuition{Intuition}\newcommand\ttlTrap{Trap}
\newcommand\ttlMethod{Method}\newcommand\ttlTask{Problem}\newcommand\ttlExample{Example}
\newcommand\ttlKey{Key result}\newcommand\ttlSrcErr{Error in the source}
\newcommand\ttlSolution{Solution.}\newcommand\ttlAnswer{Answer:}
\newcommand\studytitles[9]{%
  \renewcommand\ttlDefinition{#1}\renewcommand\ttlIntuition{#2}\renewcommand\ttlTrap{#3}%
  \renewcommand\ttlMethod{#4}\renewcommand\ttlTask{#5}\renewcommand\ttlExample{#6}%
  \renewcommand\ttlKey{#7}\renewcommand\ttlSrcErr{#8}\renewcommand\ttlSolution{#9}}
\ifdefstring{\studylang}{russian}{\studytitles{Определение}{Интуиция}{Ловушка}{Алгоритм}{Задача}{Пример}{Главное}{Ошибка в исходнике}{Решение.}\renewcommand\ttlAnswer{Ответ:}}{}
\ifdefstring{\studylang}{ngerman}{\studytitles{Definition}{Intuition}{Falle}{Vorgehen}{Aufgabe}{Beispiel}{Merke}{Fehler in der Vorlage}{Lösung.}\renewcommand\ttlAnswer{Ergebnis:}}{}
\ifdefstring{\studylang}{german}{\studytitles{Definition}{Intuition}{Falle}{Vorgehen}{Aufgabe}{Beispiel}{Merke}{Fehler in der Vorlage}{Lösung.}\renewcommand\ttlAnswer{Ergebnis:}}{}
\ifdefstring{\studylang}{french}{\studytitles{Définition}{Intuition}{Piège}{Méthode}{Exercice}{Exemple}{À retenir}{Erreur dans la source}{Solution.}\renewcommand\ttlAnswer{Réponse :}}{}
\ifdefstring{\studylang}{spanish}{\studytitles{Definición}{Intuición}{Trampa}{Método}{Ejercicio}{Ejemplo}{Clave}{Error en la fuente}{Solución.}\renewcommand\ttlAnswer{Respuesta:}}{}

% --- boxes ---
\tcbset{
  base/.style={enhanced,boxrule=0.6pt,arc=2pt,left=6pt,right=6pt,top=4pt,bottom=4pt,
    fonttitle=\bfseries\sffamily\small,before skip=6pt,after skip=6pt},
}
\NewTColorBox{defn}{O{}O{}}{base,colback=accent!4,colframe=accent!70,title={\ttlDefinition#1},#2}
\NewTColorBox{idea}{O{}O{}}{base,colback=idea!5,colframe=idea!75,title={\ttlIntuition#1},#2}
\NewTColorBox{trap}{O{}O{}}{base,colback=trap!5,colframe=trap!80,title={\ttlTrap#1},#2}
\NewTColorBox{method}{O{}O{}}{base,colback=idea!3,colframe=idea!60,title={\ttlMethod#1},#2}
\NewTColorBox{task}{O{}O{}}{base,colback=ex!5,colframe=ex!75,title={\ttlTask#1},#2}
\NewTColorBox{exmpl}{O{}O{}}{base,colback=ex!3,colframe=ex!55,title={\ttlExample#1},#2}
\NewTColorBox{keyres}{O{}O{}}{base,colback=accent!8,colframe=accent,title={\ttlKey#1},#2}
\NewTColorBox{srcerr}{O{}O{}}{base,colback=trap!3,colframe=trap!60,title={\ttlSrcErr#1},#2}
\newcommand{\sol}{\par\smallskip\noindent\textbf{\ttlSolution}\ }
\newcommand{\ans}[1]{\par\smallskip\noindent\colorbox{accent!10}{\parbox{\dimexpr\linewidth-2\fboxsep}{\textbf{\ttlAnswer} #1}}}

% --- headings (never stranded at a page bottom) ---
\titleformat{\section}{\Large\bfseries\sffamily\color{accent}}{\thesection}{0.6em}{}
\titleformat{\subsection}{\large\bfseries\sffamily\color{accent!85!black}}{\thesubsection}{0.5em}{}
\titleformat{\subsubsection}{\normalsize\bfseries\sffamily}{}{0em}{}
\titlespacing*{\section}{0pt}{14pt}{6pt}
\titlespacing*{\subsection}{0pt}{10pt}{4pt}
\titlespacing*{\subsubsection}{0pt}{8pt}{2pt}
% A heading never ends a page: it needs room for itself plus the start of what follows. Before a
% tall unbreakable box, add \needspace{<box height>} yourself or the box moves and leaves the heading.
\pretocmd{\section}{\needspace{8\baselineskip}}{}{}
\pretocmd{\subsection}{\needspace{7\baselineskip}}{}{}
\pretocmd{\subsubsection}{\needspace{5\baselineskip}}{}{}

% --- graphs ---
% \begin{axis}[econ, xlabel=$q$, ylabel=$p$]. Label curves at their end:
%   \addplot[c1,domain=0:8]{2*x+4} node[pos=1,anchor=west,font=\footnotesize]{$MC$};
% With several curves, use pos=0.9/0.95/1 and anchor=south west/west/north west so labels don't collide.
\pgfplotsset{
  every axis/.append style={tick label style={font=\footnotesize},label style={font=\small},
    legend style={font=\footnotesize,draw=none,fill=none}},
  econ/.style={axis lines=left,xlabel style={at={(ticklabel* cs:1)},anchor=north west},
    ylabel style={rotate=-90,at={(ticklabel* cs:1)},anchor=south east},
    xtick=\empty,ytick=\empty,width=6.2cm,height=5cm,samples=120,clip=false,
    every axis plot/.append style={thick},
    axis line style={-{Stealth[length=5pt]}},enlargelimits=false},
}

% Running header: \studyheader{Course · Chapter}{Deliverable}
\newcommand{\studyheader}[2]{%
  \pagestyle{fancy}\fancyhf{}%
  \fancyhead[L]{\small\sffamily\color{grayline}#1}%
  \fancyhead[R]{\small\sffamily\color{grayline}#2}%
  \fancyfoot[C]{\small\thepage}%
  \renewcommand{\headrulewidth}{0.3pt}}

% Box titles are built in for english, russian, ngerman/german, french, spanish. For another
% language, after % study-pack shared preamble (pdfLaTeX). Copy next to your .tex; don't edit it per document.
% Usage, after \documentclass and geometry:
%   \def\studylang{english}    % output language: english, russian, ngerman, french, spanish, ...
%   \def\sourcelang{english}   % language of the course materials (for \term{...}); optional
%   \input{preamble}
% Box titles are built in for english, russian, ngerman/german, french, spanish. For another
% language, after \input{preamble}, \renewcommand the \ttl... macros below.
%
% Boxes: defn idea trap method task exmpl keyres srcerr.
%   \begin{trap}[: suffix] ... \end{trap}             title "Trap: suffix"
%   \begin{exmpl}[ 3][breakable] ... \end{exmpl}      long box allowed to split across pages
% Boxes do NOT split across pages by default; pass [breakable] as the 2nd option only for a box
% longer than about half a page.
\providecommand\studylang{english}
\providecommand\sourcelang{english}
\usepackage{etoolbox}
\usepackage[T1,T2A]{fontenc}
\usepackage[utf8]{inputenc}
% Spanish babel's active shorthands (and decimal-dot handling) break TikZ, \% spacing and icomma.
\ifboolexpr{test {\ifdefstring{\studylang}{spanish}} or test {\ifdefstring{\sourcelang}{spanish}}}
  {\PassOptionsToPackage{es-noshorthands,es-nodecimaldot}{babel}}{}
\ifdefstring{\studylang}{\sourcelang}
  {\usepackage[\studylang]{babel}}
  {\usepackage[\sourcelang,\studylang]{babel}}
\usepackage{paratype}
\usepackage{amsmath,amssymb,mathtools}
% Decimal comma for non-English output. In math, type a space after any comma that separates
% items: (1, 2), (x, y), a_{1, 4}. Otherwise (1,2) prints as the decimal 1,2.
\ifdefstring{\studylang}{english}{}{\usepackage{icomma}}
\usepackage{xcolor}
\usepackage[most]{tcolorbox}
\usepackage{enumitem}
\usepackage{booktabs,array,tabularx,multirow}
\usepackage{pgfplots}
\pgfplotsset{compat=1.18}
\usetikzlibrary{arrows.meta,positioning,calc,shapes.geometric,fit,backgrounds,babel}% babel: French/Turkish active : ; ! stay safe in TikZ
\usepackage{titlesec}
\usepackage{fancyhdr}
\usepackage{multicol}
\usepackage{needspace}
\usepackage[hidelinks]{hyperref}

% --- layout defaults that prevent the usual fix-up rounds ---
\setlength{\parindent}{0pt}
\setlength{\parskip}{3pt plus 1pt}
\setlength{\emergencystretch}{3em}   % no overfull text lines
\allowdisplaybreaks                   % long align blocks may break across pages
\clubpenalty=10000 \widowpenalty=10000 \displaywidowpenalty=10000
\raggedbottom                         % no stretched pages before an unbreakable box
\setlist{nosep,leftmargin=1.4em}

\definecolor{accent}{RGB}{22,70,135}
\definecolor{idea}{RGB}{0,120,110}
\definecolor{trap}{RGB}{190,80,0}
\definecolor{ex}{RGB}{60,110,30}
\definecolor{grayline}{RGB}{150,150,150}
\definecolor{c1}{RGB}{22,70,135}
\definecolor{c2}{RGB}{200,90,20}
\definecolor{c3}{RGB}{40,140,60}
\definecolor{c4}{RGB}{170,40,90}

% Source-language term inside output-language text: \term{deadweight loss}
\newcommand{\term}[1]{\textit{\foreignlanguage{\sourcelang}{#1}}}

% --- box titles ---
\newcommand\ttlDefinition{Definition}\newcommand\ttlIntuition{Intuition}\newcommand\ttlTrap{Trap}
\newcommand\ttlMethod{Method}\newcommand\ttlTask{Problem}\newcommand\ttlExample{Example}
\newcommand\ttlKey{Key result}\newcommand\ttlSrcErr{Error in the source}
\newcommand\ttlSolution{Solution.}\newcommand\ttlAnswer{Answer:}
\newcommand\studytitles[9]{%
  \renewcommand\ttlDefinition{#1}\renewcommand\ttlIntuition{#2}\renewcommand\ttlTrap{#3}%
  \renewcommand\ttlMethod{#4}\renewcommand\ttlTask{#5}\renewcommand\ttlExample{#6}%
  \renewcommand\ttlKey{#7}\renewcommand\ttlSrcErr{#8}\renewcommand\ttlSolution{#9}}
\ifdefstring{\studylang}{russian}{\studytitles{Определение}{Интуиция}{Ловушка}{Алгоритм}{Задача}{Пример}{Главное}{Ошибка в исходнике}{Решение.}\renewcommand\ttlAnswer{Ответ:}}{}
\ifdefstring{\studylang}{ngerman}{\studytitles{Definition}{Intuition}{Falle}{Vorgehen}{Aufgabe}{Beispiel}{Merke}{Fehler in der Vorlage}{Lösung.}\renewcommand\ttlAnswer{Ergebnis:}}{}
\ifdefstring{\studylang}{german}{\studytitles{Definition}{Intuition}{Falle}{Vorgehen}{Aufgabe}{Beispiel}{Merke}{Fehler in der Vorlage}{Lösung.}\renewcommand\ttlAnswer{Ergebnis:}}{}
\ifdefstring{\studylang}{french}{\studytitles{Définition}{Intuition}{Piège}{Méthode}{Exercice}{Exemple}{À retenir}{Erreur dans la source}{Solution.}\renewcommand\ttlAnswer{Réponse :}}{}
\ifdefstring{\studylang}{spanish}{\studytitles{Definición}{Intuición}{Trampa}{Método}{Ejercicio}{Ejemplo}{Clave}{Error en la fuente}{Solución.}\renewcommand\ttlAnswer{Respuesta:}}{}

% --- boxes ---
\tcbset{
  base/.style={enhanced,boxrule=0.6pt,arc=2pt,left=6pt,right=6pt,top=4pt,bottom=4pt,
    fonttitle=\bfseries\sffamily\small,before skip=6pt,after skip=6pt},
}
\NewTColorBox{defn}{O{}O{}}{base,colback=accent!4,colframe=accent!70,title={\ttlDefinition#1},#2}
\NewTColorBox{idea}{O{}O{}}{base,colback=idea!5,colframe=idea!75,title={\ttlIntuition#1},#2}
\NewTColorBox{trap}{O{}O{}}{base,colback=trap!5,colframe=trap!80,title={\ttlTrap#1},#2}
\NewTColorBox{method}{O{}O{}}{base,colback=idea!3,colframe=idea!60,title={\ttlMethod#1},#2}
\NewTColorBox{task}{O{}O{}}{base,colback=ex!5,colframe=ex!75,title={\ttlTask#1},#2}
\NewTColorBox{exmpl}{O{}O{}}{base,colback=ex!3,colframe=ex!55,title={\ttlExample#1},#2}
\NewTColorBox{keyres}{O{}O{}}{base,colback=accent!8,colframe=accent,title={\ttlKey#1},#2}
\NewTColorBox{srcerr}{O{}O{}}{base,colback=trap!3,colframe=trap!60,title={\ttlSrcErr#1},#2}
\newcommand{\sol}{\par\smallskip\noindent\textbf{\ttlSolution}\ }
\newcommand{\ans}[1]{\par\smallskip\noindent\colorbox{accent!10}{\parbox{\dimexpr\linewidth-2\fboxsep}{\textbf{\ttlAnswer} #1}}}

% --- headings (never stranded at a page bottom) ---
\titleformat{\section}{\Large\bfseries\sffamily\color{accent}}{\thesection}{0.6em}{}
\titleformat{\subsection}{\large\bfseries\sffamily\color{accent!85!black}}{\thesubsection}{0.5em}{}
\titleformat{\subsubsection}{\normalsize\bfseries\sffamily}{}{0em}{}
\titlespacing*{\section}{0pt}{14pt}{6pt}
\titlespacing*{\subsection}{0pt}{10pt}{4pt}
\titlespacing*{\subsubsection}{0pt}{8pt}{2pt}
% A heading never ends a page: it needs room for itself plus the start of what follows. Before a
% tall unbreakable box, add \needspace{<box height>} yourself or the box moves and leaves the heading.
\pretocmd{\section}{\needspace{8\baselineskip}}{}{}
\pretocmd{\subsection}{\needspace{7\baselineskip}}{}{}
\pretocmd{\subsubsection}{\needspace{5\baselineskip}}{}{}

% --- graphs ---
% \begin{axis}[econ, xlabel=$q$, ylabel=$p$]. Label curves at their end:
%   \addplot[c1,domain=0:8]{2*x+4} node[pos=1,anchor=west,font=\footnotesize]{$MC$};
% With several curves, use pos=0.9/0.95/1 and anchor=south west/west/north west so labels don't collide.
\pgfplotsset{
  every axis/.append style={tick label style={font=\footnotesize},label style={font=\small},
    legend style={font=\footnotesize,draw=none,fill=none}},
  econ/.style={axis lines=left,xlabel style={at={(ticklabel* cs:1)},anchor=north west},
    ylabel style={rotate=-90,at={(ticklabel* cs:1)},anchor=south east},
    xtick=\empty,ytick=\empty,width=6.2cm,height=5cm,samples=120,clip=false,
    every axis plot/.append style={thick},
    axis line style={-{Stealth[length=5pt]}},enlargelimits=false},
}

% Running header: \studyheader{Course · Chapter}{Deliverable}
\newcommand{\studyheader}[2]{%
  \pagestyle{fancy}\fancyhf{}%
  \fancyhead[L]{\small\sffamily\color{grayline}#1}%
  \fancyhead[R]{\small\sffamily\color{grayline}#2}%
  \fancyfoot[C]{\small\thepage}%
  \renewcommand{\headrulewidth}{0.3pt}}
, \renewcommand the \ttl... macros below.
%
% Boxes: defn idea trap method task exmpl keyres srcerr.
%   \begin{trap}[: suffix] ... \end{trap}             title "Trap: suffix"
%   \begin{exmpl}[ 3][breakable] ... \end{exmpl}      long box allowed to split across pages
% Boxes do NOT split across pages by default; pass [breakable] as the 2nd option only for a box
% longer than about half a page.
\providecommand\studylang{english}
\providecommand\sourcelang{english}
\usepackage{etoolbox}
\usepackage[T1,T2A]{fontenc}
\usepackage[utf8]{inputenc}
% Spanish babel's active shorthands (and decimal-dot handling) break TikZ, \% spacing and icomma.
\ifboolexpr{test {\ifdefstring{\studylang}{spanish}} or test {\ifdefstring{\sourcelang}{spanish}}}
  {\PassOptionsToPackage{es-noshorthands,es-nodecimaldot}{babel}}{}
\ifdefstring{\studylang}{\sourcelang}
  {\usepackage[\studylang]{babel}}
  {\usepackage[\sourcelang,\studylang]{babel}}
\usepackage{paratype}
\usepackage{amsmath,amssymb,mathtools}
% Decimal comma for non-English output. In math, type a space after any comma that separates
% items: (1, 2), (x, y), a_{1, 4}. Otherwise (1,2) prints as the decimal 1,2.
\ifdefstring{\studylang}{english}{}{\usepackage{icomma}}
\usepackage{xcolor}
\usepackage[most]{tcolorbox}
\usepackage{enumitem}
\usepackage{booktabs,array,tabularx,multirow}
\usepackage{pgfplots}
\pgfplotsset{compat=1.18}
\usetikzlibrary{arrows.meta,positioning,calc,shapes.geometric,fit,backgrounds,babel}% babel: French/Turkish active : ; ! stay safe in TikZ
\usepackage{titlesec}
\usepackage{fancyhdr}
\usepackage{multicol}
\usepackage{needspace}
\usepackage[hidelinks]{hyperref}

% --- layout defaults that prevent the usual fix-up rounds ---
\setlength{\parindent}{0pt}
\setlength{\parskip}{3pt plus 1pt}
\setlength{\emergencystretch}{3em}   % no overfull text lines
\allowdisplaybreaks                   % long align blocks may break across pages
\clubpenalty=10000 \widowpenalty=10000 \displaywidowpenalty=10000
\raggedbottom                         % no stretched pages before an unbreakable box
\setlist{nosep,leftmargin=1.4em}

\definecolor{accent}{RGB}{22,70,135}
\definecolor{idea}{RGB}{0,120,110}
\definecolor{trap}{RGB}{190,80,0}
\definecolor{ex}{RGB}{60,110,30}
\definecolor{grayline}{RGB}{150,150,150}
\definecolor{c1}{RGB}{22,70,135}
\definecolor{c2}{RGB}{200,90,20}
\definecolor{c3}{RGB}{40,140,60}
\definecolor{c4}{RGB}{170,40,90}

% Source-language term inside output-language text: \term{deadweight loss}
\newcommand{\term}[1]{\textit{\foreignlanguage{\sourcelang}{#1}}}

% --- box titles ---
\newcommand\ttlDefinition{Definition}\newcommand\ttlIntuition{Intuition}\newcommand\ttlTrap{Trap}
\newcommand\ttlMethod{Method}\newcommand\ttlTask{Problem}\newcommand\ttlExample{Example}
\newcommand\ttlKey{Key result}\newcommand\ttlSrcErr{Error in the source}
\newcommand\ttlSolution{Solution.}\newcommand\ttlAnswer{Answer:}
\newcommand\studytitles[9]{%
  \renewcommand\ttlDefinition{#1}\renewcommand\ttlIntuition{#2}\renewcommand\ttlTrap{#3}%
  \renewcommand\ttlMethod{#4}\renewcommand\ttlTask{#5}\renewcommand\ttlExample{#6}%
  \renewcommand\ttlKey{#7}\renewcommand\ttlSrcErr{#8}\renewcommand\ttlSolution{#9}}
\ifdefstring{\studylang}{russian}{\studytitles{Определение}{Интуиция}{Ловушка}{Алгоритм}{Задача}{Пример}{Главное}{Ошибка в исходнике}{Решение.}\renewcommand\ttlAnswer{Ответ:}}{}
\ifdefstring{\studylang}{ngerman}{\studytitles{Definition}{Intuition}{Falle}{Vorgehen}{Aufgabe}{Beispiel}{Merke}{Fehler in der Vorlage}{Lösung.}\renewcommand\ttlAnswer{Ergebnis:}}{}
\ifdefstring{\studylang}{german}{\studytitles{Definition}{Intuition}{Falle}{Vorgehen}{Aufgabe}{Beispiel}{Merke}{Fehler in der Vorlage}{Lösung.}\renewcommand\ttlAnswer{Ergebnis:}}{}
\ifdefstring{\studylang}{french}{\studytitles{Définition}{Intuition}{Piège}{Méthode}{Exercice}{Exemple}{À retenir}{Erreur dans la source}{Solution.}\renewcommand\ttlAnswer{Réponse :}}{}
\ifdefstring{\studylang}{spanish}{\studytitles{Definición}{Intuición}{Trampa}{Método}{Ejercicio}{Ejemplo}{Clave}{Error en la fuente}{Solución.}\renewcommand\ttlAnswer{Respuesta:}}{}

% --- boxes ---
\tcbset{
  base/.style={enhanced,boxrule=0.6pt,arc=2pt,left=6pt,right=6pt,top=4pt,bottom=4pt,
    fonttitle=\bfseries\sffamily\small,before skip=6pt,after skip=6pt},
}
\NewTColorBox{defn}{O{}O{}}{base,colback=accent!4,colframe=accent!70,title={\ttlDefinition#1},#2}
\NewTColorBox{idea}{O{}O{}}{base,colback=idea!5,colframe=idea!75,title={\ttlIntuition#1},#2}
\NewTColorBox{trap}{O{}O{}}{base,colback=trap!5,colframe=trap!80,title={\ttlTrap#1},#2}
\NewTColorBox{method}{O{}O{}}{base,colback=idea!3,colframe=idea!60,title={\ttlMethod#1},#2}
\NewTColorBox{task}{O{}O{}}{base,colback=ex!5,colframe=ex!75,title={\ttlTask#1},#2}
\NewTColorBox{exmpl}{O{}O{}}{base,colback=ex!3,colframe=ex!55,title={\ttlExample#1},#2}
\NewTColorBox{keyres}{O{}O{}}{base,colback=accent!8,colframe=accent,title={\ttlKey#1},#2}
\NewTColorBox{srcerr}{O{}O{}}{base,colback=trap!3,colframe=trap!60,title={\ttlSrcErr#1},#2}
\newcommand{\sol}{\par\smallskip\noindent\textbf{\ttlSolution}\ }
\newcommand{\ans}[1]{\par\smallskip\noindent\colorbox{accent!10}{\parbox{\dimexpr\linewidth-2\fboxsep}{\textbf{\ttlAnswer} #1}}}

% --- headings (never stranded at a page bottom) ---
\titleformat{\section}{\Large\bfseries\sffamily\color{accent}}{\thesection}{0.6em}{}
\titleformat{\subsection}{\large\bfseries\sffamily\color{accent!85!black}}{\thesubsection}{0.5em}{}
\titleformat{\subsubsection}{\normalsize\bfseries\sffamily}{}{0em}{}
\titlespacing*{\section}{0pt}{14pt}{6pt}
\titlespacing*{\subsection}{0pt}{10pt}{4pt}
\titlespacing*{\subsubsection}{0pt}{8pt}{2pt}
% A heading never ends a page: it needs room for itself plus the start of what follows. Before a
% tall unbreakable box, add \needspace{<box height>} yourself or the box moves and leaves the heading.
\pretocmd{\section}{\needspace{8\baselineskip}}{}{}
\pretocmd{\subsection}{\needspace{7\baselineskip}}{}{}
\pretocmd{\subsubsection}{\needspace{5\baselineskip}}{}{}

% --- graphs ---
% \begin{axis}[econ, xlabel=$q$, ylabel=$p$]. Label curves at their end:
%   \addplot[c1,domain=0:8]{2*x+4} node[pos=1,anchor=west,font=\footnotesize]{$MC$};
% With several curves, use pos=0.9/0.95/1 and anchor=south west/west/north west so labels don't collide.
\pgfplotsset{
  every axis/.append style={tick label style={font=\footnotesize},label style={font=\small},
    legend style={font=\footnotesize,draw=none,fill=none}},
  econ/.style={axis lines=left,xlabel style={at={(ticklabel* cs:1)},anchor=north west},
    ylabel style={rotate=-90,at={(ticklabel* cs:1)},anchor=south east},
    xtick=\empty,ytick=\empty,width=6.2cm,height=5cm,samples=120,clip=false,
    every axis plot/.append style={thick},
    axis line style={-{Stealth[length=5pt]}},enlargelimits=false},
}

% Running header: \studyheader{Course · Chapter}{Deliverable}
\newcommand{\studyheader}[2]{%
  \pagestyle{fancy}\fancyhf{}%
  \fancyhead[L]{\small\sffamily\color{grayline}#1}%
  \fancyhead[R]{\small\sffamily\color{grayline}#2}%
  \fancyfoot[C]{\small\thepage}%
  \renewcommand{\headrulewidth}{0.3pt}}

% Box titles are built in for english, russian, ngerman/german, french, spanish. For another
% language, after % study-pack shared preamble (pdfLaTeX). Copy next to your .tex; don't edit it per document.
% Usage, after \documentclass and geometry:
%   \def\studylang{english}    % output language: english, russian, ngerman, french, spanish, ...
%   \def\sourcelang{english}   % language of the course materials (for \term{...}); optional
%   % study-pack shared preamble (pdfLaTeX). Copy next to your .tex; don't edit it per document.
% Usage, after \documentclass and geometry:
%   \def\studylang{english}    % output language: english, russian, ngerman, french, spanish, ...
%   \def\sourcelang{english}   % language of the course materials (for \term{...}); optional
%   \input{preamble}
% Box titles are built in for english, russian, ngerman/german, french, spanish. For another
% language, after \input{preamble}, \renewcommand the \ttl... macros below.
%
% Boxes: defn idea trap method task exmpl keyres srcerr.
%   \begin{trap}[: suffix] ... \end{trap}             title "Trap: suffix"
%   \begin{exmpl}[ 3][breakable] ... \end{exmpl}      long box allowed to split across pages
% Boxes do NOT split across pages by default; pass [breakable] as the 2nd option only for a box
% longer than about half a page.
\providecommand\studylang{english}
\providecommand\sourcelang{english}
\usepackage{etoolbox}
\usepackage[T1,T2A]{fontenc}
\usepackage[utf8]{inputenc}
% Spanish babel's active shorthands (and decimal-dot handling) break TikZ, \% spacing and icomma.
\ifboolexpr{test {\ifdefstring{\studylang}{spanish}} or test {\ifdefstring{\sourcelang}{spanish}}}
  {\PassOptionsToPackage{es-noshorthands,es-nodecimaldot}{babel}}{}
\ifdefstring{\studylang}{\sourcelang}
  {\usepackage[\studylang]{babel}}
  {\usepackage[\sourcelang,\studylang]{babel}}
\usepackage{paratype}
\usepackage{amsmath,amssymb,mathtools}
% Decimal comma for non-English output. In math, type a space after any comma that separates
% items: (1, 2), (x, y), a_{1, 4}. Otherwise (1,2) prints as the decimal 1,2.
\ifdefstring{\studylang}{english}{}{\usepackage{icomma}}
\usepackage{xcolor}
\usepackage[most]{tcolorbox}
\usepackage{enumitem}
\usepackage{booktabs,array,tabularx,multirow}
\usepackage{pgfplots}
\pgfplotsset{compat=1.18}
\usetikzlibrary{arrows.meta,positioning,calc,shapes.geometric,fit,backgrounds,babel}% babel: French/Turkish active : ; ! stay safe in TikZ
\usepackage{titlesec}
\usepackage{fancyhdr}
\usepackage{multicol}
\usepackage{needspace}
\usepackage[hidelinks]{hyperref}

% --- layout defaults that prevent the usual fix-up rounds ---
\setlength{\parindent}{0pt}
\setlength{\parskip}{3pt plus 1pt}
\setlength{\emergencystretch}{3em}   % no overfull text lines
\allowdisplaybreaks                   % long align blocks may break across pages
\clubpenalty=10000 \widowpenalty=10000 \displaywidowpenalty=10000
\raggedbottom                         % no stretched pages before an unbreakable box
\setlist{nosep,leftmargin=1.4em}

\definecolor{accent}{RGB}{22,70,135}
\definecolor{idea}{RGB}{0,120,110}
\definecolor{trap}{RGB}{190,80,0}
\definecolor{ex}{RGB}{60,110,30}
\definecolor{grayline}{RGB}{150,150,150}
\definecolor{c1}{RGB}{22,70,135}
\definecolor{c2}{RGB}{200,90,20}
\definecolor{c3}{RGB}{40,140,60}
\definecolor{c4}{RGB}{170,40,90}

% Source-language term inside output-language text: \term{deadweight loss}
\newcommand{\term}[1]{\textit{\foreignlanguage{\sourcelang}{#1}}}

% --- box titles ---
\newcommand\ttlDefinition{Definition}\newcommand\ttlIntuition{Intuition}\newcommand\ttlTrap{Trap}
\newcommand\ttlMethod{Method}\newcommand\ttlTask{Problem}\newcommand\ttlExample{Example}
\newcommand\ttlKey{Key result}\newcommand\ttlSrcErr{Error in the source}
\newcommand\ttlSolution{Solution.}\newcommand\ttlAnswer{Answer:}
\newcommand\studytitles[9]{%
  \renewcommand\ttlDefinition{#1}\renewcommand\ttlIntuition{#2}\renewcommand\ttlTrap{#3}%
  \renewcommand\ttlMethod{#4}\renewcommand\ttlTask{#5}\renewcommand\ttlExample{#6}%
  \renewcommand\ttlKey{#7}\renewcommand\ttlSrcErr{#8}\renewcommand\ttlSolution{#9}}
\ifdefstring{\studylang}{russian}{\studytitles{Определение}{Интуиция}{Ловушка}{Алгоритм}{Задача}{Пример}{Главное}{Ошибка в исходнике}{Решение.}\renewcommand\ttlAnswer{Ответ:}}{}
\ifdefstring{\studylang}{ngerman}{\studytitles{Definition}{Intuition}{Falle}{Vorgehen}{Aufgabe}{Beispiel}{Merke}{Fehler in der Vorlage}{Lösung.}\renewcommand\ttlAnswer{Ergebnis:}}{}
\ifdefstring{\studylang}{german}{\studytitles{Definition}{Intuition}{Falle}{Vorgehen}{Aufgabe}{Beispiel}{Merke}{Fehler in der Vorlage}{Lösung.}\renewcommand\ttlAnswer{Ergebnis:}}{}
\ifdefstring{\studylang}{french}{\studytitles{Définition}{Intuition}{Piège}{Méthode}{Exercice}{Exemple}{À retenir}{Erreur dans la source}{Solution.}\renewcommand\ttlAnswer{Réponse :}}{}
\ifdefstring{\studylang}{spanish}{\studytitles{Definición}{Intuición}{Trampa}{Método}{Ejercicio}{Ejemplo}{Clave}{Error en la fuente}{Solución.}\renewcommand\ttlAnswer{Respuesta:}}{}

% --- boxes ---
\tcbset{
  base/.style={enhanced,boxrule=0.6pt,arc=2pt,left=6pt,right=6pt,top=4pt,bottom=4pt,
    fonttitle=\bfseries\sffamily\small,before skip=6pt,after skip=6pt},
}
\NewTColorBox{defn}{O{}O{}}{base,colback=accent!4,colframe=accent!70,title={\ttlDefinition#1},#2}
\NewTColorBox{idea}{O{}O{}}{base,colback=idea!5,colframe=idea!75,title={\ttlIntuition#1},#2}
\NewTColorBox{trap}{O{}O{}}{base,colback=trap!5,colframe=trap!80,title={\ttlTrap#1},#2}
\NewTColorBox{method}{O{}O{}}{base,colback=idea!3,colframe=idea!60,title={\ttlMethod#1},#2}
\NewTColorBox{task}{O{}O{}}{base,colback=ex!5,colframe=ex!75,title={\ttlTask#1},#2}
\NewTColorBox{exmpl}{O{}O{}}{base,colback=ex!3,colframe=ex!55,title={\ttlExample#1},#2}
\NewTColorBox{keyres}{O{}O{}}{base,colback=accent!8,colframe=accent,title={\ttlKey#1},#2}
\NewTColorBox{srcerr}{O{}O{}}{base,colback=trap!3,colframe=trap!60,title={\ttlSrcErr#1},#2}
\newcommand{\sol}{\par\smallskip\noindent\textbf{\ttlSolution}\ }
\newcommand{\ans}[1]{\par\smallskip\noindent\colorbox{accent!10}{\parbox{\dimexpr\linewidth-2\fboxsep}{\textbf{\ttlAnswer} #1}}}

% --- headings (never stranded at a page bottom) ---
\titleformat{\section}{\Large\bfseries\sffamily\color{accent}}{\thesection}{0.6em}{}
\titleformat{\subsection}{\large\bfseries\sffamily\color{accent!85!black}}{\thesubsection}{0.5em}{}
\titleformat{\subsubsection}{\normalsize\bfseries\sffamily}{}{0em}{}
\titlespacing*{\section}{0pt}{14pt}{6pt}
\titlespacing*{\subsection}{0pt}{10pt}{4pt}
\titlespacing*{\subsubsection}{0pt}{8pt}{2pt}
% A heading never ends a page: it needs room for itself plus the start of what follows. Before a
% tall unbreakable box, add \needspace{<box height>} yourself or the box moves and leaves the heading.
\pretocmd{\section}{\needspace{8\baselineskip}}{}{}
\pretocmd{\subsection}{\needspace{7\baselineskip}}{}{}
\pretocmd{\subsubsection}{\needspace{5\baselineskip}}{}{}

% --- graphs ---
% \begin{axis}[econ, xlabel=$q$, ylabel=$p$]. Label curves at their end:
%   \addplot[c1,domain=0:8]{2*x+4} node[pos=1,anchor=west,font=\footnotesize]{$MC$};
% With several curves, use pos=0.9/0.95/1 and anchor=south west/west/north west so labels don't collide.
\pgfplotsset{
  every axis/.append style={tick label style={font=\footnotesize},label style={font=\small},
    legend style={font=\footnotesize,draw=none,fill=none}},
  econ/.style={axis lines=left,xlabel style={at={(ticklabel* cs:1)},anchor=north west},
    ylabel style={rotate=-90,at={(ticklabel* cs:1)},anchor=south east},
    xtick=\empty,ytick=\empty,width=6.2cm,height=5cm,samples=120,clip=false,
    every axis plot/.append style={thick},
    axis line style={-{Stealth[length=5pt]}},enlargelimits=false},
}

% Running header: \studyheader{Course · Chapter}{Deliverable}
\newcommand{\studyheader}[2]{%
  \pagestyle{fancy}\fancyhf{}%
  \fancyhead[L]{\small\sffamily\color{grayline}#1}%
  \fancyhead[R]{\small\sffamily\color{grayline}#2}%
  \fancyfoot[C]{\small\thepage}%
  \renewcommand{\headrulewidth}{0.3pt}}

% Box titles are built in for english, russian, ngerman/german, french, spanish. For another
% language, after % study-pack shared preamble (pdfLaTeX). Copy next to your .tex; don't edit it per document.
% Usage, after \documentclass and geometry:
%   \def\studylang{english}    % output language: english, russian, ngerman, french, spanish, ...
%   \def\sourcelang{english}   % language of the course materials (for \term{...}); optional
%   \input{preamble}
% Box titles are built in for english, russian, ngerman/german, french, spanish. For another
% language, after \input{preamble}, \renewcommand the \ttl... macros below.
%
% Boxes: defn idea trap method task exmpl keyres srcerr.
%   \begin{trap}[: suffix] ... \end{trap}             title "Trap: suffix"
%   \begin{exmpl}[ 3][breakable] ... \end{exmpl}      long box allowed to split across pages
% Boxes do NOT split across pages by default; pass [breakable] as the 2nd option only for a box
% longer than about half a page.
\providecommand\studylang{english}
\providecommand\sourcelang{english}
\usepackage{etoolbox}
\usepackage[T1,T2A]{fontenc}
\usepackage[utf8]{inputenc}
% Spanish babel's active shorthands (and decimal-dot handling) break TikZ, \% spacing and icomma.
\ifboolexpr{test {\ifdefstring{\studylang}{spanish}} or test {\ifdefstring{\sourcelang}{spanish}}}
  {\PassOptionsToPackage{es-noshorthands,es-nodecimaldot}{babel}}{}
\ifdefstring{\studylang}{\sourcelang}
  {\usepackage[\studylang]{babel}}
  {\usepackage[\sourcelang,\studylang]{babel}}
\usepackage{paratype}
\usepackage{amsmath,amssymb,mathtools}
% Decimal comma for non-English output. In math, type a space after any comma that separates
% items: (1, 2), (x, y), a_{1, 4}. Otherwise (1,2) prints as the decimal 1,2.
\ifdefstring{\studylang}{english}{}{\usepackage{icomma}}
\usepackage{xcolor}
\usepackage[most]{tcolorbox}
\usepackage{enumitem}
\usepackage{booktabs,array,tabularx,multirow}
\usepackage{pgfplots}
\pgfplotsset{compat=1.18}
\usetikzlibrary{arrows.meta,positioning,calc,shapes.geometric,fit,backgrounds,babel}% babel: French/Turkish active : ; ! stay safe in TikZ
\usepackage{titlesec}
\usepackage{fancyhdr}
\usepackage{multicol}
\usepackage{needspace}
\usepackage[hidelinks]{hyperref}

% --- layout defaults that prevent the usual fix-up rounds ---
\setlength{\parindent}{0pt}
\setlength{\parskip}{3pt plus 1pt}
\setlength{\emergencystretch}{3em}   % no overfull text lines
\allowdisplaybreaks                   % long align blocks may break across pages
\clubpenalty=10000 \widowpenalty=10000 \displaywidowpenalty=10000
\raggedbottom                         % no stretched pages before an unbreakable box
\setlist{nosep,leftmargin=1.4em}

\definecolor{accent}{RGB}{22,70,135}
\definecolor{idea}{RGB}{0,120,110}
\definecolor{trap}{RGB}{190,80,0}
\definecolor{ex}{RGB}{60,110,30}
\definecolor{grayline}{RGB}{150,150,150}
\definecolor{c1}{RGB}{22,70,135}
\definecolor{c2}{RGB}{200,90,20}
\definecolor{c3}{RGB}{40,140,60}
\definecolor{c4}{RGB}{170,40,90}

% Source-language term inside output-language text: \term{deadweight loss}
\newcommand{\term}[1]{\textit{\foreignlanguage{\sourcelang}{#1}}}

% --- box titles ---
\newcommand\ttlDefinition{Definition}\newcommand\ttlIntuition{Intuition}\newcommand\ttlTrap{Trap}
\newcommand\ttlMethod{Method}\newcommand\ttlTask{Problem}\newcommand\ttlExample{Example}
\newcommand\ttlKey{Key result}\newcommand\ttlSrcErr{Error in the source}
\newcommand\ttlSolution{Solution.}\newcommand\ttlAnswer{Answer:}
\newcommand\studytitles[9]{%
  \renewcommand\ttlDefinition{#1}\renewcommand\ttlIntuition{#2}\renewcommand\ttlTrap{#3}%
  \renewcommand\ttlMethod{#4}\renewcommand\ttlTask{#5}\renewcommand\ttlExample{#6}%
  \renewcommand\ttlKey{#7}\renewcommand\ttlSrcErr{#8}\renewcommand\ttlSolution{#9}}
\ifdefstring{\studylang}{russian}{\studytitles{Определение}{Интуиция}{Ловушка}{Алгоритм}{Задача}{Пример}{Главное}{Ошибка в исходнике}{Решение.}\renewcommand\ttlAnswer{Ответ:}}{}
\ifdefstring{\studylang}{ngerman}{\studytitles{Definition}{Intuition}{Falle}{Vorgehen}{Aufgabe}{Beispiel}{Merke}{Fehler in der Vorlage}{Lösung.}\renewcommand\ttlAnswer{Ergebnis:}}{}
\ifdefstring{\studylang}{german}{\studytitles{Definition}{Intuition}{Falle}{Vorgehen}{Aufgabe}{Beispiel}{Merke}{Fehler in der Vorlage}{Lösung.}\renewcommand\ttlAnswer{Ergebnis:}}{}
\ifdefstring{\studylang}{french}{\studytitles{Définition}{Intuition}{Piège}{Méthode}{Exercice}{Exemple}{À retenir}{Erreur dans la source}{Solution.}\renewcommand\ttlAnswer{Réponse :}}{}
\ifdefstring{\studylang}{spanish}{\studytitles{Definición}{Intuición}{Trampa}{Método}{Ejercicio}{Ejemplo}{Clave}{Error en la fuente}{Solución.}\renewcommand\ttlAnswer{Respuesta:}}{}

% --- boxes ---
\tcbset{
  base/.style={enhanced,boxrule=0.6pt,arc=2pt,left=6pt,right=6pt,top=4pt,bottom=4pt,
    fonttitle=\bfseries\sffamily\small,before skip=6pt,after skip=6pt},
}
\NewTColorBox{defn}{O{}O{}}{base,colback=accent!4,colframe=accent!70,title={\ttlDefinition#1},#2}
\NewTColorBox{idea}{O{}O{}}{base,colback=idea!5,colframe=idea!75,title={\ttlIntuition#1},#2}
\NewTColorBox{trap}{O{}O{}}{base,colback=trap!5,colframe=trap!80,title={\ttlTrap#1},#2}
\NewTColorBox{method}{O{}O{}}{base,colback=idea!3,colframe=idea!60,title={\ttlMethod#1},#2}
\NewTColorBox{task}{O{}O{}}{base,colback=ex!5,colframe=ex!75,title={\ttlTask#1},#2}
\NewTColorBox{exmpl}{O{}O{}}{base,colback=ex!3,colframe=ex!55,title={\ttlExample#1},#2}
\NewTColorBox{keyres}{O{}O{}}{base,colback=accent!8,colframe=accent,title={\ttlKey#1},#2}
\NewTColorBox{srcerr}{O{}O{}}{base,colback=trap!3,colframe=trap!60,title={\ttlSrcErr#1},#2}
\newcommand{\sol}{\par\smallskip\noindent\textbf{\ttlSolution}\ }
\newcommand{\ans}[1]{\par\smallskip\noindent\colorbox{accent!10}{\parbox{\dimexpr\linewidth-2\fboxsep}{\textbf{\ttlAnswer} #1}}}

% --- headings (never stranded at a page bottom) ---
\titleformat{\section}{\Large\bfseries\sffamily\color{accent}}{\thesection}{0.6em}{}
\titleformat{\subsection}{\large\bfseries\sffamily\color{accent!85!black}}{\thesubsection}{0.5em}{}
\titleformat{\subsubsection}{\normalsize\bfseries\sffamily}{}{0em}{}
\titlespacing*{\section}{0pt}{14pt}{6pt}
\titlespacing*{\subsection}{0pt}{10pt}{4pt}
\titlespacing*{\subsubsection}{0pt}{8pt}{2pt}
% A heading never ends a page: it needs room for itself plus the start of what follows. Before a
% tall unbreakable box, add \needspace{<box height>} yourself or the box moves and leaves the heading.
\pretocmd{\section}{\needspace{8\baselineskip}}{}{}
\pretocmd{\subsection}{\needspace{7\baselineskip}}{}{}
\pretocmd{\subsubsection}{\needspace{5\baselineskip}}{}{}

% --- graphs ---
% \begin{axis}[econ, xlabel=$q$, ylabel=$p$]. Label curves at their end:
%   \addplot[c1,domain=0:8]{2*x+4} node[pos=1,anchor=west,font=\footnotesize]{$MC$};
% With several curves, use pos=0.9/0.95/1 and anchor=south west/west/north west so labels don't collide.
\pgfplotsset{
  every axis/.append style={tick label style={font=\footnotesize},label style={font=\small},
    legend style={font=\footnotesize,draw=none,fill=none}},
  econ/.style={axis lines=left,xlabel style={at={(ticklabel* cs:1)},anchor=north west},
    ylabel style={rotate=-90,at={(ticklabel* cs:1)},anchor=south east},
    xtick=\empty,ytick=\empty,width=6.2cm,height=5cm,samples=120,clip=false,
    every axis plot/.append style={thick},
    axis line style={-{Stealth[length=5pt]}},enlargelimits=false},
}

% Running header: \studyheader{Course · Chapter}{Deliverable}
\newcommand{\studyheader}[2]{%
  \pagestyle{fancy}\fancyhf{}%
  \fancyhead[L]{\small\sffamily\color{grayline}#1}%
  \fancyhead[R]{\small\sffamily\color{grayline}#2}%
  \fancyfoot[C]{\small\thepage}%
  \renewcommand{\headrulewidth}{0.3pt}}
, \renewcommand the \ttl... macros below.
%
% Boxes: defn idea trap method task exmpl keyres srcerr.
%   \begin{trap}[: suffix] ... \end{trap}             title "Trap: suffix"
%   \begin{exmpl}[ 3][breakable] ... \end{exmpl}      long box allowed to split across pages
% Boxes do NOT split across pages by default; pass [breakable] as the 2nd option only for a box
% longer than about half a page.
\providecommand\studylang{english}
\providecommand\sourcelang{english}
\usepackage{etoolbox}
\usepackage[T1,T2A]{fontenc}
\usepackage[utf8]{inputenc}
% Spanish babel's active shorthands (and decimal-dot handling) break TikZ, \% spacing and icomma.
\ifboolexpr{test {\ifdefstring{\studylang}{spanish}} or test {\ifdefstring{\sourcelang}{spanish}}}
  {\PassOptionsToPackage{es-noshorthands,es-nodecimaldot}{babel}}{}
\ifdefstring{\studylang}{\sourcelang}
  {\usepackage[\studylang]{babel}}
  {\usepackage[\sourcelang,\studylang]{babel}}
\usepackage{paratype}
\usepackage{amsmath,amssymb,mathtools}
% Decimal comma for non-English output. In math, type a space after any comma that separates
% items: (1, 2), (x, y), a_{1, 4}. Otherwise (1,2) prints as the decimal 1,2.
\ifdefstring{\studylang}{english}{}{\usepackage{icomma}}
\usepackage{xcolor}
\usepackage[most]{tcolorbox}
\usepackage{enumitem}
\usepackage{booktabs,array,tabularx,multirow}
\usepackage{pgfplots}
\pgfplotsset{compat=1.18}
\usetikzlibrary{arrows.meta,positioning,calc,shapes.geometric,fit,backgrounds,babel}% babel: French/Turkish active : ; ! stay safe in TikZ
\usepackage{titlesec}
\usepackage{fancyhdr}
\usepackage{multicol}
\usepackage{needspace}
\usepackage[hidelinks]{hyperref}

% --- layout defaults that prevent the usual fix-up rounds ---
\setlength{\parindent}{0pt}
\setlength{\parskip}{3pt plus 1pt}
\setlength{\emergencystretch}{3em}   % no overfull text lines
\allowdisplaybreaks                   % long align blocks may break across pages
\clubpenalty=10000 \widowpenalty=10000 \displaywidowpenalty=10000
\raggedbottom                         % no stretched pages before an unbreakable box
\setlist{nosep,leftmargin=1.4em}

\definecolor{accent}{RGB}{22,70,135}
\definecolor{idea}{RGB}{0,120,110}
\definecolor{trap}{RGB}{190,80,0}
\definecolor{ex}{RGB}{60,110,30}
\definecolor{grayline}{RGB}{150,150,150}
\definecolor{c1}{RGB}{22,70,135}
\definecolor{c2}{RGB}{200,90,20}
\definecolor{c3}{RGB}{40,140,60}
\definecolor{c4}{RGB}{170,40,90}

% Source-language term inside output-language text: \term{deadweight loss}
\newcommand{\term}[1]{\textit{\foreignlanguage{\sourcelang}{#1}}}

% --- box titles ---
\newcommand\ttlDefinition{Definition}\newcommand\ttlIntuition{Intuition}\newcommand\ttlTrap{Trap}
\newcommand\ttlMethod{Method}\newcommand\ttlTask{Problem}\newcommand\ttlExample{Example}
\newcommand\ttlKey{Key result}\newcommand\ttlSrcErr{Error in the source}
\newcommand\ttlSolution{Solution.}\newcommand\ttlAnswer{Answer:}
\newcommand\studytitles[9]{%
  \renewcommand\ttlDefinition{#1}\renewcommand\ttlIntuition{#2}\renewcommand\ttlTrap{#3}%
  \renewcommand\ttlMethod{#4}\renewcommand\ttlTask{#5}\renewcommand\ttlExample{#6}%
  \renewcommand\ttlKey{#7}\renewcommand\ttlSrcErr{#8}\renewcommand\ttlSolution{#9}}
\ifdefstring{\studylang}{russian}{\studytitles{Определение}{Интуиция}{Ловушка}{Алгоритм}{Задача}{Пример}{Главное}{Ошибка в исходнике}{Решение.}\renewcommand\ttlAnswer{Ответ:}}{}
\ifdefstring{\studylang}{ngerman}{\studytitles{Definition}{Intuition}{Falle}{Vorgehen}{Aufgabe}{Beispiel}{Merke}{Fehler in der Vorlage}{Lösung.}\renewcommand\ttlAnswer{Ergebnis:}}{}
\ifdefstring{\studylang}{german}{\studytitles{Definition}{Intuition}{Falle}{Vorgehen}{Aufgabe}{Beispiel}{Merke}{Fehler in der Vorlage}{Lösung.}\renewcommand\ttlAnswer{Ergebnis:}}{}
\ifdefstring{\studylang}{french}{\studytitles{Définition}{Intuition}{Piège}{Méthode}{Exercice}{Exemple}{À retenir}{Erreur dans la source}{Solution.}\renewcommand\ttlAnswer{Réponse :}}{}
\ifdefstring{\studylang}{spanish}{\studytitles{Definición}{Intuición}{Trampa}{Método}{Ejercicio}{Ejemplo}{Clave}{Error en la fuente}{Solución.}\renewcommand\ttlAnswer{Respuesta:}}{}

% --- boxes ---
\tcbset{
  base/.style={enhanced,boxrule=0.6pt,arc=2pt,left=6pt,right=6pt,top=4pt,bottom=4pt,
    fonttitle=\bfseries\sffamily\small,before skip=6pt,after skip=6pt},
}
\NewTColorBox{defn}{O{}O{}}{base,colback=accent!4,colframe=accent!70,title={\ttlDefinition#1},#2}
\NewTColorBox{idea}{O{}O{}}{base,colback=idea!5,colframe=idea!75,title={\ttlIntuition#1},#2}
\NewTColorBox{trap}{O{}O{}}{base,colback=trap!5,colframe=trap!80,title={\ttlTrap#1},#2}
\NewTColorBox{method}{O{}O{}}{base,colback=idea!3,colframe=idea!60,title={\ttlMethod#1},#2}
\NewTColorBox{task}{O{}O{}}{base,colback=ex!5,colframe=ex!75,title={\ttlTask#1},#2}
\NewTColorBox{exmpl}{O{}O{}}{base,colback=ex!3,colframe=ex!55,title={\ttlExample#1},#2}
\NewTColorBox{keyres}{O{}O{}}{base,colback=accent!8,colframe=accent,title={\ttlKey#1},#2}
\NewTColorBox{srcerr}{O{}O{}}{base,colback=trap!3,colframe=trap!60,title={\ttlSrcErr#1},#2}
\newcommand{\sol}{\par\smallskip\noindent\textbf{\ttlSolution}\ }
\newcommand{\ans}[1]{\par\smallskip\noindent\colorbox{accent!10}{\parbox{\dimexpr\linewidth-2\fboxsep}{\textbf{\ttlAnswer} #1}}}

% --- headings (never stranded at a page bottom) ---
\titleformat{\section}{\Large\bfseries\sffamily\color{accent}}{\thesection}{0.6em}{}
\titleformat{\subsection}{\large\bfseries\sffamily\color{accent!85!black}}{\thesubsection}{0.5em}{}
\titleformat{\subsubsection}{\normalsize\bfseries\sffamily}{}{0em}{}
\titlespacing*{\section}{0pt}{14pt}{6pt}
\titlespacing*{\subsection}{0pt}{10pt}{4pt}
\titlespacing*{\subsubsection}{0pt}{8pt}{2pt}
% A heading never ends a page: it needs room for itself plus the start of what follows. Before a
% tall unbreakable box, add \needspace{<box height>} yourself or the box moves and leaves the heading.
\pretocmd{\section}{\needspace{8\baselineskip}}{}{}
\pretocmd{\subsection}{\needspace{7\baselineskip}}{}{}
\pretocmd{\subsubsection}{\needspace{5\baselineskip}}{}{}

% --- graphs ---
% \begin{axis}[econ, xlabel=$q$, ylabel=$p$]. Label curves at their end:
%   \addplot[c1,domain=0:8]{2*x+4} node[pos=1,anchor=west,font=\footnotesize]{$MC$};
% With several curves, use pos=0.9/0.95/1 and anchor=south west/west/north west so labels don't collide.
\pgfplotsset{
  every axis/.append style={tick label style={font=\footnotesize},label style={font=\small},
    legend style={font=\footnotesize,draw=none,fill=none}},
  econ/.style={axis lines=left,xlabel style={at={(ticklabel* cs:1)},anchor=north west},
    ylabel style={rotate=-90,at={(ticklabel* cs:1)},anchor=south east},
    xtick=\empty,ytick=\empty,width=6.2cm,height=5cm,samples=120,clip=false,
    every axis plot/.append style={thick},
    axis line style={-{Stealth[length=5pt]}},enlargelimits=false},
}

% Running header: \studyheader{Course · Chapter}{Deliverable}
\newcommand{\studyheader}[2]{%
  \pagestyle{fancy}\fancyhf{}%
  \fancyhead[L]{\small\sffamily\color{grayline}#1}%
  \fancyhead[R]{\small\sffamily\color{grayline}#2}%
  \fancyfoot[C]{\small\thepage}%
  \renewcommand{\headrulewidth}{0.3pt}}
, \renewcommand the \ttl... macros below.
%
% Boxes: defn idea trap method task exmpl keyres srcerr.
%   \begin{trap}[: suffix] ... \end{trap}             title "Trap: suffix"
%   \begin{exmpl}[ 3][breakable] ... \end{exmpl}      long box allowed to split across pages
% Boxes do NOT split across pages by default; pass [breakable] as the 2nd option only for a box
% longer than about half a page.
\providecommand\studylang{english}
\providecommand\sourcelang{english}
\usepackage{etoolbox}
\usepackage[T1,T2A]{fontenc}
\usepackage[utf8]{inputenc}
% Spanish babel's active shorthands (and decimal-dot handling) break TikZ, \% spacing and icomma.
\ifboolexpr{test {\ifdefstring{\studylang}{spanish}} or test {\ifdefstring{\sourcelang}{spanish}}}
  {\PassOptionsToPackage{es-noshorthands,es-nodecimaldot}{babel}}{}
\ifdefstring{\studylang}{\sourcelang}
  {\usepackage[\studylang]{babel}}
  {\usepackage[\sourcelang,\studylang]{babel}}
\usepackage{paratype}
\usepackage{amsmath,amssymb,mathtools}
% Decimal comma for non-English output. In math, type a space after any comma that separates
% items: (1, 2), (x, y), a_{1, 4}. Otherwise (1,2) prints as the decimal 1,2.
\ifdefstring{\studylang}{english}{}{\usepackage{icomma}}
\usepackage{xcolor}
\usepackage[most]{tcolorbox}
\usepackage{enumitem}
\usepackage{booktabs,array,tabularx,multirow}
\usepackage{pgfplots}
\pgfplotsset{compat=1.18}
\usetikzlibrary{arrows.meta,positioning,calc,shapes.geometric,fit,backgrounds,babel}% babel: French/Turkish active : ; ! stay safe in TikZ
\usepackage{titlesec}
\usepackage{fancyhdr}
\usepackage{multicol}
\usepackage{needspace}
\usepackage[hidelinks]{hyperref}

% --- layout defaults that prevent the usual fix-up rounds ---
\setlength{\parindent}{0pt}
\setlength{\parskip}{3pt plus 1pt}
\setlength{\emergencystretch}{3em}   % no overfull text lines
\allowdisplaybreaks                   % long align blocks may break across pages
\clubpenalty=10000 \widowpenalty=10000 \displaywidowpenalty=10000
\raggedbottom                         % no stretched pages before an unbreakable box
\setlist{nosep,leftmargin=1.4em}

\definecolor{accent}{RGB}{22,70,135}
\definecolor{idea}{RGB}{0,120,110}
\definecolor{trap}{RGB}{190,80,0}
\definecolor{ex}{RGB}{60,110,30}
\definecolor{grayline}{RGB}{150,150,150}
\definecolor{c1}{RGB}{22,70,135}
\definecolor{c2}{RGB}{200,90,20}
\definecolor{c3}{RGB}{40,140,60}
\definecolor{c4}{RGB}{170,40,90}

% Source-language term inside output-language text: \term{deadweight loss}
\newcommand{\term}[1]{\textit{\foreignlanguage{\sourcelang}{#1}}}

% --- box titles ---
\newcommand\ttlDefinition{Definition}\newcommand\ttlIntuition{Intuition}\newcommand\ttlTrap{Trap}
\newcommand\ttlMethod{Method}\newcommand\ttlTask{Problem}\newcommand\ttlExample{Example}
\newcommand\ttlKey{Key result}\newcommand\ttlSrcErr{Error in the source}
\newcommand\ttlSolution{Solution.}\newcommand\ttlAnswer{Answer:}
\newcommand\studytitles[9]{%
  \renewcommand\ttlDefinition{#1}\renewcommand\ttlIntuition{#2}\renewcommand\ttlTrap{#3}%
  \renewcommand\ttlMethod{#4}\renewcommand\ttlTask{#5}\renewcommand\ttlExample{#6}%
  \renewcommand\ttlKey{#7}\renewcommand\ttlSrcErr{#8}\renewcommand\ttlSolution{#9}}
\ifdefstring{\studylang}{russian}{\studytitles{Определение}{Интуиция}{Ловушка}{Алгоритм}{Задача}{Пример}{Главное}{Ошибка в исходнике}{Решение.}\renewcommand\ttlAnswer{Ответ:}}{}
\ifdefstring{\studylang}{ngerman}{\studytitles{Definition}{Intuition}{Falle}{Vorgehen}{Aufgabe}{Beispiel}{Merke}{Fehler in der Vorlage}{Lösung.}\renewcommand\ttlAnswer{Ergebnis:}}{}
\ifdefstring{\studylang}{german}{\studytitles{Definition}{Intuition}{Falle}{Vorgehen}{Aufgabe}{Beispiel}{Merke}{Fehler in der Vorlage}{Lösung.}\renewcommand\ttlAnswer{Ergebnis:}}{}
\ifdefstring{\studylang}{french}{\studytitles{Définition}{Intuition}{Piège}{Méthode}{Exercice}{Exemple}{À retenir}{Erreur dans la source}{Solution.}\renewcommand\ttlAnswer{Réponse :}}{}
\ifdefstring{\studylang}{spanish}{\studytitles{Definición}{Intuición}{Trampa}{Método}{Ejercicio}{Ejemplo}{Clave}{Error en la fuente}{Solución.}\renewcommand\ttlAnswer{Respuesta:}}{}

% --- boxes ---
\tcbset{
  base/.style={enhanced,boxrule=0.6pt,arc=2pt,left=6pt,right=6pt,top=4pt,bottom=4pt,
    fonttitle=\bfseries\sffamily\small,before skip=6pt,after skip=6pt},
}
\NewTColorBox{defn}{O{}O{}}{base,colback=accent!4,colframe=accent!70,title={\ttlDefinition#1},#2}
\NewTColorBox{idea}{O{}O{}}{base,colback=idea!5,colframe=idea!75,title={\ttlIntuition#1},#2}
\NewTColorBox{trap}{O{}O{}}{base,colback=trap!5,colframe=trap!80,title={\ttlTrap#1},#2}
\NewTColorBox{method}{O{}O{}}{base,colback=idea!3,colframe=idea!60,title={\ttlMethod#1},#2}
\NewTColorBox{task}{O{}O{}}{base,colback=ex!5,colframe=ex!75,title={\ttlTask#1},#2}
\NewTColorBox{exmpl}{O{}O{}}{base,colback=ex!3,colframe=ex!55,title={\ttlExample#1},#2}
\NewTColorBox{keyres}{O{}O{}}{base,colback=accent!8,colframe=accent,title={\ttlKey#1},#2}
\NewTColorBox{srcerr}{O{}O{}}{base,colback=trap!3,colframe=trap!60,title={\ttlSrcErr#1},#2}
\newcommand{\sol}{\par\smallskip\noindent\textbf{\ttlSolution}\ }
\newcommand{\ans}[1]{\par\smallskip\noindent\colorbox{accent!10}{\parbox{\dimexpr\linewidth-2\fboxsep}{\textbf{\ttlAnswer} #1}}}

% --- headings (never stranded at a page bottom) ---
\titleformat{\section}{\Large\bfseries\sffamily\color{accent}}{\thesection}{0.6em}{}
\titleformat{\subsection}{\large\bfseries\sffamily\color{accent!85!black}}{\thesubsection}{0.5em}{}
\titleformat{\subsubsection}{\normalsize\bfseries\sffamily}{}{0em}{}
\titlespacing*{\section}{0pt}{14pt}{6pt}
\titlespacing*{\subsection}{0pt}{10pt}{4pt}
\titlespacing*{\subsubsection}{0pt}{8pt}{2pt}
% A heading never ends a page: it needs room for itself plus the start of what follows. Before a
% tall unbreakable box, add \needspace{<box height>} yourself or the box moves and leaves the heading.
\pretocmd{\section}{\needspace{8\baselineskip}}{}{}
\pretocmd{\subsection}{\needspace{7\baselineskip}}{}{}
\pretocmd{\subsubsection}{\needspace{5\baselineskip}}{}{}

% --- graphs ---
% \begin{axis}[econ, xlabel=$q$, ylabel=$p$]. Label curves at their end:
%   \addplot[c1,domain=0:8]{2*x+4} node[pos=1,anchor=west,font=\footnotesize]{$MC$};
% With several curves, use pos=0.9/0.95/1 and anchor=south west/west/north west so labels don't collide.
\pgfplotsset{
  every axis/.append style={tick label style={font=\footnotesize},label style={font=\small},
    legend style={font=\footnotesize,draw=none,fill=none}},
  econ/.style={axis lines=left,xlabel style={at={(ticklabel* cs:1)},anchor=north west},
    ylabel style={rotate=-90,at={(ticklabel* cs:1)},anchor=south east},
    xtick=\empty,ytick=\empty,width=6.2cm,height=5cm,samples=120,clip=false,
    every axis plot/.append style={thick},
    axis line style={-{Stealth[length=5pt]}},enlargelimits=false},
}

% Running header: \studyheader{Course · Chapter}{Deliverable}
\newcommand{\studyheader}[2]{%
  \pagestyle{fancy}\fancyhf{}%
  \fancyhead[L]{\small\sffamily\color{grayline}#1}%
  \fancyhead[R]{\small\sffamily\color{grayline}#2}%
  \fancyfoot[C]{\small\thepage}%
  \renewcommand{\headrulewidth}{0.3pt}}

% Box titles are built in for english, russian, ngerman/german, french, spanish. For another
% language, after % study-pack shared preamble (pdfLaTeX). Copy next to your .tex; don't edit it per document.
% Usage, after \documentclass and geometry:
%   \def\studylang{english}    % output language: english, russian, ngerman, french, spanish, ...
%   \def\sourcelang{english}   % language of the course materials (for \term{...}); optional
%   % study-pack shared preamble (pdfLaTeX). Copy next to your .tex; don't edit it per document.
% Usage, after \documentclass and geometry:
%   \def\studylang{english}    % output language: english, russian, ngerman, french, spanish, ...
%   \def\sourcelang{english}   % language of the course materials (for \term{...}); optional
%   % study-pack shared preamble (pdfLaTeX). Copy next to your .tex; don't edit it per document.
% Usage, after \documentclass and geometry:
%   \def\studylang{english}    % output language: english, russian, ngerman, french, spanish, ...
%   \def\sourcelang{english}   % language of the course materials (for \term{...}); optional
%   \input{preamble}
% Box titles are built in for english, russian, ngerman/german, french, spanish. For another
% language, after \input{preamble}, \renewcommand the \ttl... macros below.
%
% Boxes: defn idea trap method task exmpl keyres srcerr.
%   \begin{trap}[: suffix] ... \end{trap}             title "Trap: suffix"
%   \begin{exmpl}[ 3][breakable] ... \end{exmpl}      long box allowed to split across pages
% Boxes do NOT split across pages by default; pass [breakable] as the 2nd option only for a box
% longer than about half a page.
\providecommand\studylang{english}
\providecommand\sourcelang{english}
\usepackage{etoolbox}
\usepackage[T1,T2A]{fontenc}
\usepackage[utf8]{inputenc}
% Spanish babel's active shorthands (and decimal-dot handling) break TikZ, \% spacing and icomma.
\ifboolexpr{test {\ifdefstring{\studylang}{spanish}} or test {\ifdefstring{\sourcelang}{spanish}}}
  {\PassOptionsToPackage{es-noshorthands,es-nodecimaldot}{babel}}{}
\ifdefstring{\studylang}{\sourcelang}
  {\usepackage[\studylang]{babel}}
  {\usepackage[\sourcelang,\studylang]{babel}}
\usepackage{paratype}
\usepackage{amsmath,amssymb,mathtools}
% Decimal comma for non-English output. In math, type a space after any comma that separates
% items: (1, 2), (x, y), a_{1, 4}. Otherwise (1,2) prints as the decimal 1,2.
\ifdefstring{\studylang}{english}{}{\usepackage{icomma}}
\usepackage{xcolor}
\usepackage[most]{tcolorbox}
\usepackage{enumitem}
\usepackage{booktabs,array,tabularx,multirow}
\usepackage{pgfplots}
\pgfplotsset{compat=1.18}
\usetikzlibrary{arrows.meta,positioning,calc,shapes.geometric,fit,backgrounds,babel}% babel: French/Turkish active : ; ! stay safe in TikZ
\usepackage{titlesec}
\usepackage{fancyhdr}
\usepackage{multicol}
\usepackage{needspace}
\usepackage[hidelinks]{hyperref}

% --- layout defaults that prevent the usual fix-up rounds ---
\setlength{\parindent}{0pt}
\setlength{\parskip}{3pt plus 1pt}
\setlength{\emergencystretch}{3em}   % no overfull text lines
\allowdisplaybreaks                   % long align blocks may break across pages
\clubpenalty=10000 \widowpenalty=10000 \displaywidowpenalty=10000
\raggedbottom                         % no stretched pages before an unbreakable box
\setlist{nosep,leftmargin=1.4em}

\definecolor{accent}{RGB}{22,70,135}
\definecolor{idea}{RGB}{0,120,110}
\definecolor{trap}{RGB}{190,80,0}
\definecolor{ex}{RGB}{60,110,30}
\definecolor{grayline}{RGB}{150,150,150}
\definecolor{c1}{RGB}{22,70,135}
\definecolor{c2}{RGB}{200,90,20}
\definecolor{c3}{RGB}{40,140,60}
\definecolor{c4}{RGB}{170,40,90}

% Source-language term inside output-language text: \term{deadweight loss}
\newcommand{\term}[1]{\textit{\foreignlanguage{\sourcelang}{#1}}}

% --- box titles ---
\newcommand\ttlDefinition{Definition}\newcommand\ttlIntuition{Intuition}\newcommand\ttlTrap{Trap}
\newcommand\ttlMethod{Method}\newcommand\ttlTask{Problem}\newcommand\ttlExample{Example}
\newcommand\ttlKey{Key result}\newcommand\ttlSrcErr{Error in the source}
\newcommand\ttlSolution{Solution.}\newcommand\ttlAnswer{Answer:}
\newcommand\studytitles[9]{%
  \renewcommand\ttlDefinition{#1}\renewcommand\ttlIntuition{#2}\renewcommand\ttlTrap{#3}%
  \renewcommand\ttlMethod{#4}\renewcommand\ttlTask{#5}\renewcommand\ttlExample{#6}%
  \renewcommand\ttlKey{#7}\renewcommand\ttlSrcErr{#8}\renewcommand\ttlSolution{#9}}
\ifdefstring{\studylang}{russian}{\studytitles{Определение}{Интуиция}{Ловушка}{Алгоритм}{Задача}{Пример}{Главное}{Ошибка в исходнике}{Решение.}\renewcommand\ttlAnswer{Ответ:}}{}
\ifdefstring{\studylang}{ngerman}{\studytitles{Definition}{Intuition}{Falle}{Vorgehen}{Aufgabe}{Beispiel}{Merke}{Fehler in der Vorlage}{Lösung.}\renewcommand\ttlAnswer{Ergebnis:}}{}
\ifdefstring{\studylang}{german}{\studytitles{Definition}{Intuition}{Falle}{Vorgehen}{Aufgabe}{Beispiel}{Merke}{Fehler in der Vorlage}{Lösung.}\renewcommand\ttlAnswer{Ergebnis:}}{}
\ifdefstring{\studylang}{french}{\studytitles{Définition}{Intuition}{Piège}{Méthode}{Exercice}{Exemple}{À retenir}{Erreur dans la source}{Solution.}\renewcommand\ttlAnswer{Réponse :}}{}
\ifdefstring{\studylang}{spanish}{\studytitles{Definición}{Intuición}{Trampa}{Método}{Ejercicio}{Ejemplo}{Clave}{Error en la fuente}{Solución.}\renewcommand\ttlAnswer{Respuesta:}}{}

% --- boxes ---
\tcbset{
  base/.style={enhanced,boxrule=0.6pt,arc=2pt,left=6pt,right=6pt,top=4pt,bottom=4pt,
    fonttitle=\bfseries\sffamily\small,before skip=6pt,after skip=6pt},
}
\NewTColorBox{defn}{O{}O{}}{base,colback=accent!4,colframe=accent!70,title={\ttlDefinition#1},#2}
\NewTColorBox{idea}{O{}O{}}{base,colback=idea!5,colframe=idea!75,title={\ttlIntuition#1},#2}
\NewTColorBox{trap}{O{}O{}}{base,colback=trap!5,colframe=trap!80,title={\ttlTrap#1},#2}
\NewTColorBox{method}{O{}O{}}{base,colback=idea!3,colframe=idea!60,title={\ttlMethod#1},#2}
\NewTColorBox{task}{O{}O{}}{base,colback=ex!5,colframe=ex!75,title={\ttlTask#1},#2}
\NewTColorBox{exmpl}{O{}O{}}{base,colback=ex!3,colframe=ex!55,title={\ttlExample#1},#2}
\NewTColorBox{keyres}{O{}O{}}{base,colback=accent!8,colframe=accent,title={\ttlKey#1},#2}
\NewTColorBox{srcerr}{O{}O{}}{base,colback=trap!3,colframe=trap!60,title={\ttlSrcErr#1},#2}
\newcommand{\sol}{\par\smallskip\noindent\textbf{\ttlSolution}\ }
\newcommand{\ans}[1]{\par\smallskip\noindent\colorbox{accent!10}{\parbox{\dimexpr\linewidth-2\fboxsep}{\textbf{\ttlAnswer} #1}}}

% --- headings (never stranded at a page bottom) ---
\titleformat{\section}{\Large\bfseries\sffamily\color{accent}}{\thesection}{0.6em}{}
\titleformat{\subsection}{\large\bfseries\sffamily\color{accent!85!black}}{\thesubsection}{0.5em}{}
\titleformat{\subsubsection}{\normalsize\bfseries\sffamily}{}{0em}{}
\titlespacing*{\section}{0pt}{14pt}{6pt}
\titlespacing*{\subsection}{0pt}{10pt}{4pt}
\titlespacing*{\subsubsection}{0pt}{8pt}{2pt}
% A heading never ends a page: it needs room for itself plus the start of what follows. Before a
% tall unbreakable box, add \needspace{<box height>} yourself or the box moves and leaves the heading.
\pretocmd{\section}{\needspace{8\baselineskip}}{}{}
\pretocmd{\subsection}{\needspace{7\baselineskip}}{}{}
\pretocmd{\subsubsection}{\needspace{5\baselineskip}}{}{}

% --- graphs ---
% \begin{axis}[econ, xlabel=$q$, ylabel=$p$]. Label curves at their end:
%   \addplot[c1,domain=0:8]{2*x+4} node[pos=1,anchor=west,font=\footnotesize]{$MC$};
% With several curves, use pos=0.9/0.95/1 and anchor=south west/west/north west so labels don't collide.
\pgfplotsset{
  every axis/.append style={tick label style={font=\footnotesize},label style={font=\small},
    legend style={font=\footnotesize,draw=none,fill=none}},
  econ/.style={axis lines=left,xlabel style={at={(ticklabel* cs:1)},anchor=north west},
    ylabel style={rotate=-90,at={(ticklabel* cs:1)},anchor=south east},
    xtick=\empty,ytick=\empty,width=6.2cm,height=5cm,samples=120,clip=false,
    every axis plot/.append style={thick},
    axis line style={-{Stealth[length=5pt]}},enlargelimits=false},
}

% Running header: \studyheader{Course · Chapter}{Deliverable}
\newcommand{\studyheader}[2]{%
  \pagestyle{fancy}\fancyhf{}%
  \fancyhead[L]{\small\sffamily\color{grayline}#1}%
  \fancyhead[R]{\small\sffamily\color{grayline}#2}%
  \fancyfoot[C]{\small\thepage}%
  \renewcommand{\headrulewidth}{0.3pt}}

% Box titles are built in for english, russian, ngerman/german, french, spanish. For another
% language, after % study-pack shared preamble (pdfLaTeX). Copy next to your .tex; don't edit it per document.
% Usage, after \documentclass and geometry:
%   \def\studylang{english}    % output language: english, russian, ngerman, french, spanish, ...
%   \def\sourcelang{english}   % language of the course materials (for \term{...}); optional
%   \input{preamble}
% Box titles are built in for english, russian, ngerman/german, french, spanish. For another
% language, after \input{preamble}, \renewcommand the \ttl... macros below.
%
% Boxes: defn idea trap method task exmpl keyres srcerr.
%   \begin{trap}[: suffix] ... \end{trap}             title "Trap: suffix"
%   \begin{exmpl}[ 3][breakable] ... \end{exmpl}      long box allowed to split across pages
% Boxes do NOT split across pages by default; pass [breakable] as the 2nd option only for a box
% longer than about half a page.
\providecommand\studylang{english}
\providecommand\sourcelang{english}
\usepackage{etoolbox}
\usepackage[T1,T2A]{fontenc}
\usepackage[utf8]{inputenc}
% Spanish babel's active shorthands (and decimal-dot handling) break TikZ, \% spacing and icomma.
\ifboolexpr{test {\ifdefstring{\studylang}{spanish}} or test {\ifdefstring{\sourcelang}{spanish}}}
  {\PassOptionsToPackage{es-noshorthands,es-nodecimaldot}{babel}}{}
\ifdefstring{\studylang}{\sourcelang}
  {\usepackage[\studylang]{babel}}
  {\usepackage[\sourcelang,\studylang]{babel}}
\usepackage{paratype}
\usepackage{amsmath,amssymb,mathtools}
% Decimal comma for non-English output. In math, type a space after any comma that separates
% items: (1, 2), (x, y), a_{1, 4}. Otherwise (1,2) prints as the decimal 1,2.
\ifdefstring{\studylang}{english}{}{\usepackage{icomma}}
\usepackage{xcolor}
\usepackage[most]{tcolorbox}
\usepackage{enumitem}
\usepackage{booktabs,array,tabularx,multirow}
\usepackage{pgfplots}
\pgfplotsset{compat=1.18}
\usetikzlibrary{arrows.meta,positioning,calc,shapes.geometric,fit,backgrounds,babel}% babel: French/Turkish active : ; ! stay safe in TikZ
\usepackage{titlesec}
\usepackage{fancyhdr}
\usepackage{multicol}
\usepackage{needspace}
\usepackage[hidelinks]{hyperref}

% --- layout defaults that prevent the usual fix-up rounds ---
\setlength{\parindent}{0pt}
\setlength{\parskip}{3pt plus 1pt}
\setlength{\emergencystretch}{3em}   % no overfull text lines
\allowdisplaybreaks                   % long align blocks may break across pages
\clubpenalty=10000 \widowpenalty=10000 \displaywidowpenalty=10000
\raggedbottom                         % no stretched pages before an unbreakable box
\setlist{nosep,leftmargin=1.4em}

\definecolor{accent}{RGB}{22,70,135}
\definecolor{idea}{RGB}{0,120,110}
\definecolor{trap}{RGB}{190,80,0}
\definecolor{ex}{RGB}{60,110,30}
\definecolor{grayline}{RGB}{150,150,150}
\definecolor{c1}{RGB}{22,70,135}
\definecolor{c2}{RGB}{200,90,20}
\definecolor{c3}{RGB}{40,140,60}
\definecolor{c4}{RGB}{170,40,90}

% Source-language term inside output-language text: \term{deadweight loss}
\newcommand{\term}[1]{\textit{\foreignlanguage{\sourcelang}{#1}}}

% --- box titles ---
\newcommand\ttlDefinition{Definition}\newcommand\ttlIntuition{Intuition}\newcommand\ttlTrap{Trap}
\newcommand\ttlMethod{Method}\newcommand\ttlTask{Problem}\newcommand\ttlExample{Example}
\newcommand\ttlKey{Key result}\newcommand\ttlSrcErr{Error in the source}
\newcommand\ttlSolution{Solution.}\newcommand\ttlAnswer{Answer:}
\newcommand\studytitles[9]{%
  \renewcommand\ttlDefinition{#1}\renewcommand\ttlIntuition{#2}\renewcommand\ttlTrap{#3}%
  \renewcommand\ttlMethod{#4}\renewcommand\ttlTask{#5}\renewcommand\ttlExample{#6}%
  \renewcommand\ttlKey{#7}\renewcommand\ttlSrcErr{#8}\renewcommand\ttlSolution{#9}}
\ifdefstring{\studylang}{russian}{\studytitles{Определение}{Интуиция}{Ловушка}{Алгоритм}{Задача}{Пример}{Главное}{Ошибка в исходнике}{Решение.}\renewcommand\ttlAnswer{Ответ:}}{}
\ifdefstring{\studylang}{ngerman}{\studytitles{Definition}{Intuition}{Falle}{Vorgehen}{Aufgabe}{Beispiel}{Merke}{Fehler in der Vorlage}{Lösung.}\renewcommand\ttlAnswer{Ergebnis:}}{}
\ifdefstring{\studylang}{german}{\studytitles{Definition}{Intuition}{Falle}{Vorgehen}{Aufgabe}{Beispiel}{Merke}{Fehler in der Vorlage}{Lösung.}\renewcommand\ttlAnswer{Ergebnis:}}{}
\ifdefstring{\studylang}{french}{\studytitles{Définition}{Intuition}{Piège}{Méthode}{Exercice}{Exemple}{À retenir}{Erreur dans la source}{Solution.}\renewcommand\ttlAnswer{Réponse :}}{}
\ifdefstring{\studylang}{spanish}{\studytitles{Definición}{Intuición}{Trampa}{Método}{Ejercicio}{Ejemplo}{Clave}{Error en la fuente}{Solución.}\renewcommand\ttlAnswer{Respuesta:}}{}

% --- boxes ---
\tcbset{
  base/.style={enhanced,boxrule=0.6pt,arc=2pt,left=6pt,right=6pt,top=4pt,bottom=4pt,
    fonttitle=\bfseries\sffamily\small,before skip=6pt,after skip=6pt},
}
\NewTColorBox{defn}{O{}O{}}{base,colback=accent!4,colframe=accent!70,title={\ttlDefinition#1},#2}
\NewTColorBox{idea}{O{}O{}}{base,colback=idea!5,colframe=idea!75,title={\ttlIntuition#1},#2}
\NewTColorBox{trap}{O{}O{}}{base,colback=trap!5,colframe=trap!80,title={\ttlTrap#1},#2}
\NewTColorBox{method}{O{}O{}}{base,colback=idea!3,colframe=idea!60,title={\ttlMethod#1},#2}
\NewTColorBox{task}{O{}O{}}{base,colback=ex!5,colframe=ex!75,title={\ttlTask#1},#2}
\NewTColorBox{exmpl}{O{}O{}}{base,colback=ex!3,colframe=ex!55,title={\ttlExample#1},#2}
\NewTColorBox{keyres}{O{}O{}}{base,colback=accent!8,colframe=accent,title={\ttlKey#1},#2}
\NewTColorBox{srcerr}{O{}O{}}{base,colback=trap!3,colframe=trap!60,title={\ttlSrcErr#1},#2}
\newcommand{\sol}{\par\smallskip\noindent\textbf{\ttlSolution}\ }
\newcommand{\ans}[1]{\par\smallskip\noindent\colorbox{accent!10}{\parbox{\dimexpr\linewidth-2\fboxsep}{\textbf{\ttlAnswer} #1}}}

% --- headings (never stranded at a page bottom) ---
\titleformat{\section}{\Large\bfseries\sffamily\color{accent}}{\thesection}{0.6em}{}
\titleformat{\subsection}{\large\bfseries\sffamily\color{accent!85!black}}{\thesubsection}{0.5em}{}
\titleformat{\subsubsection}{\normalsize\bfseries\sffamily}{}{0em}{}
\titlespacing*{\section}{0pt}{14pt}{6pt}
\titlespacing*{\subsection}{0pt}{10pt}{4pt}
\titlespacing*{\subsubsection}{0pt}{8pt}{2pt}
% A heading never ends a page: it needs room for itself plus the start of what follows. Before a
% tall unbreakable box, add \needspace{<box height>} yourself or the box moves and leaves the heading.
\pretocmd{\section}{\needspace{8\baselineskip}}{}{}
\pretocmd{\subsection}{\needspace{7\baselineskip}}{}{}
\pretocmd{\subsubsection}{\needspace{5\baselineskip}}{}{}

% --- graphs ---
% \begin{axis}[econ, xlabel=$q$, ylabel=$p$]. Label curves at their end:
%   \addplot[c1,domain=0:8]{2*x+4} node[pos=1,anchor=west,font=\footnotesize]{$MC$};
% With several curves, use pos=0.9/0.95/1 and anchor=south west/west/north west so labels don't collide.
\pgfplotsset{
  every axis/.append style={tick label style={font=\footnotesize},label style={font=\small},
    legend style={font=\footnotesize,draw=none,fill=none}},
  econ/.style={axis lines=left,xlabel style={at={(ticklabel* cs:1)},anchor=north west},
    ylabel style={rotate=-90,at={(ticklabel* cs:1)},anchor=south east},
    xtick=\empty,ytick=\empty,width=6.2cm,height=5cm,samples=120,clip=false,
    every axis plot/.append style={thick},
    axis line style={-{Stealth[length=5pt]}},enlargelimits=false},
}

% Running header: \studyheader{Course · Chapter}{Deliverable}
\newcommand{\studyheader}[2]{%
  \pagestyle{fancy}\fancyhf{}%
  \fancyhead[L]{\small\sffamily\color{grayline}#1}%
  \fancyhead[R]{\small\sffamily\color{grayline}#2}%
  \fancyfoot[C]{\small\thepage}%
  \renewcommand{\headrulewidth}{0.3pt}}
, \renewcommand the \ttl... macros below.
%
% Boxes: defn idea trap method task exmpl keyres srcerr.
%   \begin{trap}[: suffix] ... \end{trap}             title "Trap: suffix"
%   \begin{exmpl}[ 3][breakable] ... \end{exmpl}      long box allowed to split across pages
% Boxes do NOT split across pages by default; pass [breakable] as the 2nd option only for a box
% longer than about half a page.
\providecommand\studylang{english}
\providecommand\sourcelang{english}
\usepackage{etoolbox}
\usepackage[T1,T2A]{fontenc}
\usepackage[utf8]{inputenc}
% Spanish babel's active shorthands (and decimal-dot handling) break TikZ, \% spacing and icomma.
\ifboolexpr{test {\ifdefstring{\studylang}{spanish}} or test {\ifdefstring{\sourcelang}{spanish}}}
  {\PassOptionsToPackage{es-noshorthands,es-nodecimaldot}{babel}}{}
\ifdefstring{\studylang}{\sourcelang}
  {\usepackage[\studylang]{babel}}
  {\usepackage[\sourcelang,\studylang]{babel}}
\usepackage{paratype}
\usepackage{amsmath,amssymb,mathtools}
% Decimal comma for non-English output. In math, type a space after any comma that separates
% items: (1, 2), (x, y), a_{1, 4}. Otherwise (1,2) prints as the decimal 1,2.
\ifdefstring{\studylang}{english}{}{\usepackage{icomma}}
\usepackage{xcolor}
\usepackage[most]{tcolorbox}
\usepackage{enumitem}
\usepackage{booktabs,array,tabularx,multirow}
\usepackage{pgfplots}
\pgfplotsset{compat=1.18}
\usetikzlibrary{arrows.meta,positioning,calc,shapes.geometric,fit,backgrounds,babel}% babel: French/Turkish active : ; ! stay safe in TikZ
\usepackage{titlesec}
\usepackage{fancyhdr}
\usepackage{multicol}
\usepackage{needspace}
\usepackage[hidelinks]{hyperref}

% --- layout defaults that prevent the usual fix-up rounds ---
\setlength{\parindent}{0pt}
\setlength{\parskip}{3pt plus 1pt}
\setlength{\emergencystretch}{3em}   % no overfull text lines
\allowdisplaybreaks                   % long align blocks may break across pages
\clubpenalty=10000 \widowpenalty=10000 \displaywidowpenalty=10000
\raggedbottom                         % no stretched pages before an unbreakable box
\setlist{nosep,leftmargin=1.4em}

\definecolor{accent}{RGB}{22,70,135}
\definecolor{idea}{RGB}{0,120,110}
\definecolor{trap}{RGB}{190,80,0}
\definecolor{ex}{RGB}{60,110,30}
\definecolor{grayline}{RGB}{150,150,150}
\definecolor{c1}{RGB}{22,70,135}
\definecolor{c2}{RGB}{200,90,20}
\definecolor{c3}{RGB}{40,140,60}
\definecolor{c4}{RGB}{170,40,90}

% Source-language term inside output-language text: \term{deadweight loss}
\newcommand{\term}[1]{\textit{\foreignlanguage{\sourcelang}{#1}}}

% --- box titles ---
\newcommand\ttlDefinition{Definition}\newcommand\ttlIntuition{Intuition}\newcommand\ttlTrap{Trap}
\newcommand\ttlMethod{Method}\newcommand\ttlTask{Problem}\newcommand\ttlExample{Example}
\newcommand\ttlKey{Key result}\newcommand\ttlSrcErr{Error in the source}
\newcommand\ttlSolution{Solution.}\newcommand\ttlAnswer{Answer:}
\newcommand\studytitles[9]{%
  \renewcommand\ttlDefinition{#1}\renewcommand\ttlIntuition{#2}\renewcommand\ttlTrap{#3}%
  \renewcommand\ttlMethod{#4}\renewcommand\ttlTask{#5}\renewcommand\ttlExample{#6}%
  \renewcommand\ttlKey{#7}\renewcommand\ttlSrcErr{#8}\renewcommand\ttlSolution{#9}}
\ifdefstring{\studylang}{russian}{\studytitles{Определение}{Интуиция}{Ловушка}{Алгоритм}{Задача}{Пример}{Главное}{Ошибка в исходнике}{Решение.}\renewcommand\ttlAnswer{Ответ:}}{}
\ifdefstring{\studylang}{ngerman}{\studytitles{Definition}{Intuition}{Falle}{Vorgehen}{Aufgabe}{Beispiel}{Merke}{Fehler in der Vorlage}{Lösung.}\renewcommand\ttlAnswer{Ergebnis:}}{}
\ifdefstring{\studylang}{german}{\studytitles{Definition}{Intuition}{Falle}{Vorgehen}{Aufgabe}{Beispiel}{Merke}{Fehler in der Vorlage}{Lösung.}\renewcommand\ttlAnswer{Ergebnis:}}{}
\ifdefstring{\studylang}{french}{\studytitles{Définition}{Intuition}{Piège}{Méthode}{Exercice}{Exemple}{À retenir}{Erreur dans la source}{Solution.}\renewcommand\ttlAnswer{Réponse :}}{}
\ifdefstring{\studylang}{spanish}{\studytitles{Definición}{Intuición}{Trampa}{Método}{Ejercicio}{Ejemplo}{Clave}{Error en la fuente}{Solución.}\renewcommand\ttlAnswer{Respuesta:}}{}

% --- boxes ---
\tcbset{
  base/.style={enhanced,boxrule=0.6pt,arc=2pt,left=6pt,right=6pt,top=4pt,bottom=4pt,
    fonttitle=\bfseries\sffamily\small,before skip=6pt,after skip=6pt},
}
\NewTColorBox{defn}{O{}O{}}{base,colback=accent!4,colframe=accent!70,title={\ttlDefinition#1},#2}
\NewTColorBox{idea}{O{}O{}}{base,colback=idea!5,colframe=idea!75,title={\ttlIntuition#1},#2}
\NewTColorBox{trap}{O{}O{}}{base,colback=trap!5,colframe=trap!80,title={\ttlTrap#1},#2}
\NewTColorBox{method}{O{}O{}}{base,colback=idea!3,colframe=idea!60,title={\ttlMethod#1},#2}
\NewTColorBox{task}{O{}O{}}{base,colback=ex!5,colframe=ex!75,title={\ttlTask#1},#2}
\NewTColorBox{exmpl}{O{}O{}}{base,colback=ex!3,colframe=ex!55,title={\ttlExample#1},#2}
\NewTColorBox{keyres}{O{}O{}}{base,colback=accent!8,colframe=accent,title={\ttlKey#1},#2}
\NewTColorBox{srcerr}{O{}O{}}{base,colback=trap!3,colframe=trap!60,title={\ttlSrcErr#1},#2}
\newcommand{\sol}{\par\smallskip\noindent\textbf{\ttlSolution}\ }
\newcommand{\ans}[1]{\par\smallskip\noindent\colorbox{accent!10}{\parbox{\dimexpr\linewidth-2\fboxsep}{\textbf{\ttlAnswer} #1}}}

% --- headings (never stranded at a page bottom) ---
\titleformat{\section}{\Large\bfseries\sffamily\color{accent}}{\thesection}{0.6em}{}
\titleformat{\subsection}{\large\bfseries\sffamily\color{accent!85!black}}{\thesubsection}{0.5em}{}
\titleformat{\subsubsection}{\normalsize\bfseries\sffamily}{}{0em}{}
\titlespacing*{\section}{0pt}{14pt}{6pt}
\titlespacing*{\subsection}{0pt}{10pt}{4pt}
\titlespacing*{\subsubsection}{0pt}{8pt}{2pt}
% A heading never ends a page: it needs room for itself plus the start of what follows. Before a
% tall unbreakable box, add \needspace{<box height>} yourself or the box moves and leaves the heading.
\pretocmd{\section}{\needspace{8\baselineskip}}{}{}
\pretocmd{\subsection}{\needspace{7\baselineskip}}{}{}
\pretocmd{\subsubsection}{\needspace{5\baselineskip}}{}{}

% --- graphs ---
% \begin{axis}[econ, xlabel=$q$, ylabel=$p$]. Label curves at their end:
%   \addplot[c1,domain=0:8]{2*x+4} node[pos=1,anchor=west,font=\footnotesize]{$MC$};
% With several curves, use pos=0.9/0.95/1 and anchor=south west/west/north west so labels don't collide.
\pgfplotsset{
  every axis/.append style={tick label style={font=\footnotesize},label style={font=\small},
    legend style={font=\footnotesize,draw=none,fill=none}},
  econ/.style={axis lines=left,xlabel style={at={(ticklabel* cs:1)},anchor=north west},
    ylabel style={rotate=-90,at={(ticklabel* cs:1)},anchor=south east},
    xtick=\empty,ytick=\empty,width=6.2cm,height=5cm,samples=120,clip=false,
    every axis plot/.append style={thick},
    axis line style={-{Stealth[length=5pt]}},enlargelimits=false},
}

% Running header: \studyheader{Course · Chapter}{Deliverable}
\newcommand{\studyheader}[2]{%
  \pagestyle{fancy}\fancyhf{}%
  \fancyhead[L]{\small\sffamily\color{grayline}#1}%
  \fancyhead[R]{\small\sffamily\color{grayline}#2}%
  \fancyfoot[C]{\small\thepage}%
  \renewcommand{\headrulewidth}{0.3pt}}

% Box titles are built in for english, russian, ngerman/german, french, spanish. For another
% language, after % study-pack shared preamble (pdfLaTeX). Copy next to your .tex; don't edit it per document.
% Usage, after \documentclass and geometry:
%   \def\studylang{english}    % output language: english, russian, ngerman, french, spanish, ...
%   \def\sourcelang{english}   % language of the course materials (for \term{...}); optional
%   % study-pack shared preamble (pdfLaTeX). Copy next to your .tex; don't edit it per document.
% Usage, after \documentclass and geometry:
%   \def\studylang{english}    % output language: english, russian, ngerman, french, spanish, ...
%   \def\sourcelang{english}   % language of the course materials (for \term{...}); optional
%   \input{preamble}
% Box titles are built in for english, russian, ngerman/german, french, spanish. For another
% language, after \input{preamble}, \renewcommand the \ttl... macros below.
%
% Boxes: defn idea trap method task exmpl keyres srcerr.
%   \begin{trap}[: suffix] ... \end{trap}             title "Trap: suffix"
%   \begin{exmpl}[ 3][breakable] ... \end{exmpl}      long box allowed to split across pages
% Boxes do NOT split across pages by default; pass [breakable] as the 2nd option only for a box
% longer than about half a page.
\providecommand\studylang{english}
\providecommand\sourcelang{english}
\usepackage{etoolbox}
\usepackage[T1,T2A]{fontenc}
\usepackage[utf8]{inputenc}
% Spanish babel's active shorthands (and decimal-dot handling) break TikZ, \% spacing and icomma.
\ifboolexpr{test {\ifdefstring{\studylang}{spanish}} or test {\ifdefstring{\sourcelang}{spanish}}}
  {\PassOptionsToPackage{es-noshorthands,es-nodecimaldot}{babel}}{}
\ifdefstring{\studylang}{\sourcelang}
  {\usepackage[\studylang]{babel}}
  {\usepackage[\sourcelang,\studylang]{babel}}
\usepackage{paratype}
\usepackage{amsmath,amssymb,mathtools}
% Decimal comma for non-English output. In math, type a space after any comma that separates
% items: (1, 2), (x, y), a_{1, 4}. Otherwise (1,2) prints as the decimal 1,2.
\ifdefstring{\studylang}{english}{}{\usepackage{icomma}}
\usepackage{xcolor}
\usepackage[most]{tcolorbox}
\usepackage{enumitem}
\usepackage{booktabs,array,tabularx,multirow}
\usepackage{pgfplots}
\pgfplotsset{compat=1.18}
\usetikzlibrary{arrows.meta,positioning,calc,shapes.geometric,fit,backgrounds,babel}% babel: French/Turkish active : ; ! stay safe in TikZ
\usepackage{titlesec}
\usepackage{fancyhdr}
\usepackage{multicol}
\usepackage{needspace}
\usepackage[hidelinks]{hyperref}

% --- layout defaults that prevent the usual fix-up rounds ---
\setlength{\parindent}{0pt}
\setlength{\parskip}{3pt plus 1pt}
\setlength{\emergencystretch}{3em}   % no overfull text lines
\allowdisplaybreaks                   % long align blocks may break across pages
\clubpenalty=10000 \widowpenalty=10000 \displaywidowpenalty=10000
\raggedbottom                         % no stretched pages before an unbreakable box
\setlist{nosep,leftmargin=1.4em}

\definecolor{accent}{RGB}{22,70,135}
\definecolor{idea}{RGB}{0,120,110}
\definecolor{trap}{RGB}{190,80,0}
\definecolor{ex}{RGB}{60,110,30}
\definecolor{grayline}{RGB}{150,150,150}
\definecolor{c1}{RGB}{22,70,135}
\definecolor{c2}{RGB}{200,90,20}
\definecolor{c3}{RGB}{40,140,60}
\definecolor{c4}{RGB}{170,40,90}

% Source-language term inside output-language text: \term{deadweight loss}
\newcommand{\term}[1]{\textit{\foreignlanguage{\sourcelang}{#1}}}

% --- box titles ---
\newcommand\ttlDefinition{Definition}\newcommand\ttlIntuition{Intuition}\newcommand\ttlTrap{Trap}
\newcommand\ttlMethod{Method}\newcommand\ttlTask{Problem}\newcommand\ttlExample{Example}
\newcommand\ttlKey{Key result}\newcommand\ttlSrcErr{Error in the source}
\newcommand\ttlSolution{Solution.}\newcommand\ttlAnswer{Answer:}
\newcommand\studytitles[9]{%
  \renewcommand\ttlDefinition{#1}\renewcommand\ttlIntuition{#2}\renewcommand\ttlTrap{#3}%
  \renewcommand\ttlMethod{#4}\renewcommand\ttlTask{#5}\renewcommand\ttlExample{#6}%
  \renewcommand\ttlKey{#7}\renewcommand\ttlSrcErr{#8}\renewcommand\ttlSolution{#9}}
\ifdefstring{\studylang}{russian}{\studytitles{Определение}{Интуиция}{Ловушка}{Алгоритм}{Задача}{Пример}{Главное}{Ошибка в исходнике}{Решение.}\renewcommand\ttlAnswer{Ответ:}}{}
\ifdefstring{\studylang}{ngerman}{\studytitles{Definition}{Intuition}{Falle}{Vorgehen}{Aufgabe}{Beispiel}{Merke}{Fehler in der Vorlage}{Lösung.}\renewcommand\ttlAnswer{Ergebnis:}}{}
\ifdefstring{\studylang}{german}{\studytitles{Definition}{Intuition}{Falle}{Vorgehen}{Aufgabe}{Beispiel}{Merke}{Fehler in der Vorlage}{Lösung.}\renewcommand\ttlAnswer{Ergebnis:}}{}
\ifdefstring{\studylang}{french}{\studytitles{Définition}{Intuition}{Piège}{Méthode}{Exercice}{Exemple}{À retenir}{Erreur dans la source}{Solution.}\renewcommand\ttlAnswer{Réponse :}}{}
\ifdefstring{\studylang}{spanish}{\studytitles{Definición}{Intuición}{Trampa}{Método}{Ejercicio}{Ejemplo}{Clave}{Error en la fuente}{Solución.}\renewcommand\ttlAnswer{Respuesta:}}{}

% --- boxes ---
\tcbset{
  base/.style={enhanced,boxrule=0.6pt,arc=2pt,left=6pt,right=6pt,top=4pt,bottom=4pt,
    fonttitle=\bfseries\sffamily\small,before skip=6pt,after skip=6pt},
}
\NewTColorBox{defn}{O{}O{}}{base,colback=accent!4,colframe=accent!70,title={\ttlDefinition#1},#2}
\NewTColorBox{idea}{O{}O{}}{base,colback=idea!5,colframe=idea!75,title={\ttlIntuition#1},#2}
\NewTColorBox{trap}{O{}O{}}{base,colback=trap!5,colframe=trap!80,title={\ttlTrap#1},#2}
\NewTColorBox{method}{O{}O{}}{base,colback=idea!3,colframe=idea!60,title={\ttlMethod#1},#2}
\NewTColorBox{task}{O{}O{}}{base,colback=ex!5,colframe=ex!75,title={\ttlTask#1},#2}
\NewTColorBox{exmpl}{O{}O{}}{base,colback=ex!3,colframe=ex!55,title={\ttlExample#1},#2}
\NewTColorBox{keyres}{O{}O{}}{base,colback=accent!8,colframe=accent,title={\ttlKey#1},#2}
\NewTColorBox{srcerr}{O{}O{}}{base,colback=trap!3,colframe=trap!60,title={\ttlSrcErr#1},#2}
\newcommand{\sol}{\par\smallskip\noindent\textbf{\ttlSolution}\ }
\newcommand{\ans}[1]{\par\smallskip\noindent\colorbox{accent!10}{\parbox{\dimexpr\linewidth-2\fboxsep}{\textbf{\ttlAnswer} #1}}}

% --- headings (never stranded at a page bottom) ---
\titleformat{\section}{\Large\bfseries\sffamily\color{accent}}{\thesection}{0.6em}{}
\titleformat{\subsection}{\large\bfseries\sffamily\color{accent!85!black}}{\thesubsection}{0.5em}{}
\titleformat{\subsubsection}{\normalsize\bfseries\sffamily}{}{0em}{}
\titlespacing*{\section}{0pt}{14pt}{6pt}
\titlespacing*{\subsection}{0pt}{10pt}{4pt}
\titlespacing*{\subsubsection}{0pt}{8pt}{2pt}
% A heading never ends a page: it needs room for itself plus the start of what follows. Before a
% tall unbreakable box, add \needspace{<box height>} yourself or the box moves and leaves the heading.
\pretocmd{\section}{\needspace{8\baselineskip}}{}{}
\pretocmd{\subsection}{\needspace{7\baselineskip}}{}{}
\pretocmd{\subsubsection}{\needspace{5\baselineskip}}{}{}

% --- graphs ---
% \begin{axis}[econ, xlabel=$q$, ylabel=$p$]. Label curves at their end:
%   \addplot[c1,domain=0:8]{2*x+4} node[pos=1,anchor=west,font=\footnotesize]{$MC$};
% With several curves, use pos=0.9/0.95/1 and anchor=south west/west/north west so labels don't collide.
\pgfplotsset{
  every axis/.append style={tick label style={font=\footnotesize},label style={font=\small},
    legend style={font=\footnotesize,draw=none,fill=none}},
  econ/.style={axis lines=left,xlabel style={at={(ticklabel* cs:1)},anchor=north west},
    ylabel style={rotate=-90,at={(ticklabel* cs:1)},anchor=south east},
    xtick=\empty,ytick=\empty,width=6.2cm,height=5cm,samples=120,clip=false,
    every axis plot/.append style={thick},
    axis line style={-{Stealth[length=5pt]}},enlargelimits=false},
}

% Running header: \studyheader{Course · Chapter}{Deliverable}
\newcommand{\studyheader}[2]{%
  \pagestyle{fancy}\fancyhf{}%
  \fancyhead[L]{\small\sffamily\color{grayline}#1}%
  \fancyhead[R]{\small\sffamily\color{grayline}#2}%
  \fancyfoot[C]{\small\thepage}%
  \renewcommand{\headrulewidth}{0.3pt}}

% Box titles are built in for english, russian, ngerman/german, french, spanish. For another
% language, after % study-pack shared preamble (pdfLaTeX). Copy next to your .tex; don't edit it per document.
% Usage, after \documentclass and geometry:
%   \def\studylang{english}    % output language: english, russian, ngerman, french, spanish, ...
%   \def\sourcelang{english}   % language of the course materials (for \term{...}); optional
%   \input{preamble}
% Box titles are built in for english, russian, ngerman/german, french, spanish. For another
% language, after \input{preamble}, \renewcommand the \ttl... macros below.
%
% Boxes: defn idea trap method task exmpl keyres srcerr.
%   \begin{trap}[: suffix] ... \end{trap}             title "Trap: suffix"
%   \begin{exmpl}[ 3][breakable] ... \end{exmpl}      long box allowed to split across pages
% Boxes do NOT split across pages by default; pass [breakable] as the 2nd option only for a box
% longer than about half a page.
\providecommand\studylang{english}
\providecommand\sourcelang{english}
\usepackage{etoolbox}
\usepackage[T1,T2A]{fontenc}
\usepackage[utf8]{inputenc}
% Spanish babel's active shorthands (and decimal-dot handling) break TikZ, \% spacing and icomma.
\ifboolexpr{test {\ifdefstring{\studylang}{spanish}} or test {\ifdefstring{\sourcelang}{spanish}}}
  {\PassOptionsToPackage{es-noshorthands,es-nodecimaldot}{babel}}{}
\ifdefstring{\studylang}{\sourcelang}
  {\usepackage[\studylang]{babel}}
  {\usepackage[\sourcelang,\studylang]{babel}}
\usepackage{paratype}
\usepackage{amsmath,amssymb,mathtools}
% Decimal comma for non-English output. In math, type a space after any comma that separates
% items: (1, 2), (x, y), a_{1, 4}. Otherwise (1,2) prints as the decimal 1,2.
\ifdefstring{\studylang}{english}{}{\usepackage{icomma}}
\usepackage{xcolor}
\usepackage[most]{tcolorbox}
\usepackage{enumitem}
\usepackage{booktabs,array,tabularx,multirow}
\usepackage{pgfplots}
\pgfplotsset{compat=1.18}
\usetikzlibrary{arrows.meta,positioning,calc,shapes.geometric,fit,backgrounds,babel}% babel: French/Turkish active : ; ! stay safe in TikZ
\usepackage{titlesec}
\usepackage{fancyhdr}
\usepackage{multicol}
\usepackage{needspace}
\usepackage[hidelinks]{hyperref}

% --- layout defaults that prevent the usual fix-up rounds ---
\setlength{\parindent}{0pt}
\setlength{\parskip}{3pt plus 1pt}
\setlength{\emergencystretch}{3em}   % no overfull text lines
\allowdisplaybreaks                   % long align blocks may break across pages
\clubpenalty=10000 \widowpenalty=10000 \displaywidowpenalty=10000
\raggedbottom                         % no stretched pages before an unbreakable box
\setlist{nosep,leftmargin=1.4em}

\definecolor{accent}{RGB}{22,70,135}
\definecolor{idea}{RGB}{0,120,110}
\definecolor{trap}{RGB}{190,80,0}
\definecolor{ex}{RGB}{60,110,30}
\definecolor{grayline}{RGB}{150,150,150}
\definecolor{c1}{RGB}{22,70,135}
\definecolor{c2}{RGB}{200,90,20}
\definecolor{c3}{RGB}{40,140,60}
\definecolor{c4}{RGB}{170,40,90}

% Source-language term inside output-language text: \term{deadweight loss}
\newcommand{\term}[1]{\textit{\foreignlanguage{\sourcelang}{#1}}}

% --- box titles ---
\newcommand\ttlDefinition{Definition}\newcommand\ttlIntuition{Intuition}\newcommand\ttlTrap{Trap}
\newcommand\ttlMethod{Method}\newcommand\ttlTask{Problem}\newcommand\ttlExample{Example}
\newcommand\ttlKey{Key result}\newcommand\ttlSrcErr{Error in the source}
\newcommand\ttlSolution{Solution.}\newcommand\ttlAnswer{Answer:}
\newcommand\studytitles[9]{%
  \renewcommand\ttlDefinition{#1}\renewcommand\ttlIntuition{#2}\renewcommand\ttlTrap{#3}%
  \renewcommand\ttlMethod{#4}\renewcommand\ttlTask{#5}\renewcommand\ttlExample{#6}%
  \renewcommand\ttlKey{#7}\renewcommand\ttlSrcErr{#8}\renewcommand\ttlSolution{#9}}
\ifdefstring{\studylang}{russian}{\studytitles{Определение}{Интуиция}{Ловушка}{Алгоритм}{Задача}{Пример}{Главное}{Ошибка в исходнике}{Решение.}\renewcommand\ttlAnswer{Ответ:}}{}
\ifdefstring{\studylang}{ngerman}{\studytitles{Definition}{Intuition}{Falle}{Vorgehen}{Aufgabe}{Beispiel}{Merke}{Fehler in der Vorlage}{Lösung.}\renewcommand\ttlAnswer{Ergebnis:}}{}
\ifdefstring{\studylang}{german}{\studytitles{Definition}{Intuition}{Falle}{Vorgehen}{Aufgabe}{Beispiel}{Merke}{Fehler in der Vorlage}{Lösung.}\renewcommand\ttlAnswer{Ergebnis:}}{}
\ifdefstring{\studylang}{french}{\studytitles{Définition}{Intuition}{Piège}{Méthode}{Exercice}{Exemple}{À retenir}{Erreur dans la source}{Solution.}\renewcommand\ttlAnswer{Réponse :}}{}
\ifdefstring{\studylang}{spanish}{\studytitles{Definición}{Intuición}{Trampa}{Método}{Ejercicio}{Ejemplo}{Clave}{Error en la fuente}{Solución.}\renewcommand\ttlAnswer{Respuesta:}}{}

% --- boxes ---
\tcbset{
  base/.style={enhanced,boxrule=0.6pt,arc=2pt,left=6pt,right=6pt,top=4pt,bottom=4pt,
    fonttitle=\bfseries\sffamily\small,before skip=6pt,after skip=6pt},
}
\NewTColorBox{defn}{O{}O{}}{base,colback=accent!4,colframe=accent!70,title={\ttlDefinition#1},#2}
\NewTColorBox{idea}{O{}O{}}{base,colback=idea!5,colframe=idea!75,title={\ttlIntuition#1},#2}
\NewTColorBox{trap}{O{}O{}}{base,colback=trap!5,colframe=trap!80,title={\ttlTrap#1},#2}
\NewTColorBox{method}{O{}O{}}{base,colback=idea!3,colframe=idea!60,title={\ttlMethod#1},#2}
\NewTColorBox{task}{O{}O{}}{base,colback=ex!5,colframe=ex!75,title={\ttlTask#1},#2}
\NewTColorBox{exmpl}{O{}O{}}{base,colback=ex!3,colframe=ex!55,title={\ttlExample#1},#2}
\NewTColorBox{keyres}{O{}O{}}{base,colback=accent!8,colframe=accent,title={\ttlKey#1},#2}
\NewTColorBox{srcerr}{O{}O{}}{base,colback=trap!3,colframe=trap!60,title={\ttlSrcErr#1},#2}
\newcommand{\sol}{\par\smallskip\noindent\textbf{\ttlSolution}\ }
\newcommand{\ans}[1]{\par\smallskip\noindent\colorbox{accent!10}{\parbox{\dimexpr\linewidth-2\fboxsep}{\textbf{\ttlAnswer} #1}}}

% --- headings (never stranded at a page bottom) ---
\titleformat{\section}{\Large\bfseries\sffamily\color{accent}}{\thesection}{0.6em}{}
\titleformat{\subsection}{\large\bfseries\sffamily\color{accent!85!black}}{\thesubsection}{0.5em}{}
\titleformat{\subsubsection}{\normalsize\bfseries\sffamily}{}{0em}{}
\titlespacing*{\section}{0pt}{14pt}{6pt}
\titlespacing*{\subsection}{0pt}{10pt}{4pt}
\titlespacing*{\subsubsection}{0pt}{8pt}{2pt}
% A heading never ends a page: it needs room for itself plus the start of what follows. Before a
% tall unbreakable box, add \needspace{<box height>} yourself or the box moves and leaves the heading.
\pretocmd{\section}{\needspace{8\baselineskip}}{}{}
\pretocmd{\subsection}{\needspace{7\baselineskip}}{}{}
\pretocmd{\subsubsection}{\needspace{5\baselineskip}}{}{}

% --- graphs ---
% \begin{axis}[econ, xlabel=$q$, ylabel=$p$]. Label curves at their end:
%   \addplot[c1,domain=0:8]{2*x+4} node[pos=1,anchor=west,font=\footnotesize]{$MC$};
% With several curves, use pos=0.9/0.95/1 and anchor=south west/west/north west so labels don't collide.
\pgfplotsset{
  every axis/.append style={tick label style={font=\footnotesize},label style={font=\small},
    legend style={font=\footnotesize,draw=none,fill=none}},
  econ/.style={axis lines=left,xlabel style={at={(ticklabel* cs:1)},anchor=north west},
    ylabel style={rotate=-90,at={(ticklabel* cs:1)},anchor=south east},
    xtick=\empty,ytick=\empty,width=6.2cm,height=5cm,samples=120,clip=false,
    every axis plot/.append style={thick},
    axis line style={-{Stealth[length=5pt]}},enlargelimits=false},
}

% Running header: \studyheader{Course · Chapter}{Deliverable}
\newcommand{\studyheader}[2]{%
  \pagestyle{fancy}\fancyhf{}%
  \fancyhead[L]{\small\sffamily\color{grayline}#1}%
  \fancyhead[R]{\small\sffamily\color{grayline}#2}%
  \fancyfoot[C]{\small\thepage}%
  \renewcommand{\headrulewidth}{0.3pt}}
, \renewcommand the \ttl... macros below.
%
% Boxes: defn idea trap method task exmpl keyres srcerr.
%   \begin{trap}[: suffix] ... \end{trap}             title "Trap: suffix"
%   \begin{exmpl}[ 3][breakable] ... \end{exmpl}      long box allowed to split across pages
% Boxes do NOT split across pages by default; pass [breakable] as the 2nd option only for a box
% longer than about half a page.
\providecommand\studylang{english}
\providecommand\sourcelang{english}
\usepackage{etoolbox}
\usepackage[T1,T2A]{fontenc}
\usepackage[utf8]{inputenc}
% Spanish babel's active shorthands (and decimal-dot handling) break TikZ, \% spacing and icomma.
\ifboolexpr{test {\ifdefstring{\studylang}{spanish}} or test {\ifdefstring{\sourcelang}{spanish}}}
  {\PassOptionsToPackage{es-noshorthands,es-nodecimaldot}{babel}}{}
\ifdefstring{\studylang}{\sourcelang}
  {\usepackage[\studylang]{babel}}
  {\usepackage[\sourcelang,\studylang]{babel}}
\usepackage{paratype}
\usepackage{amsmath,amssymb,mathtools}
% Decimal comma for non-English output. In math, type a space after any comma that separates
% items: (1, 2), (x, y), a_{1, 4}. Otherwise (1,2) prints as the decimal 1,2.
\ifdefstring{\studylang}{english}{}{\usepackage{icomma}}
\usepackage{xcolor}
\usepackage[most]{tcolorbox}
\usepackage{enumitem}
\usepackage{booktabs,array,tabularx,multirow}
\usepackage{pgfplots}
\pgfplotsset{compat=1.18}
\usetikzlibrary{arrows.meta,positioning,calc,shapes.geometric,fit,backgrounds,babel}% babel: French/Turkish active : ; ! stay safe in TikZ
\usepackage{titlesec}
\usepackage{fancyhdr}
\usepackage{multicol}
\usepackage{needspace}
\usepackage[hidelinks]{hyperref}

% --- layout defaults that prevent the usual fix-up rounds ---
\setlength{\parindent}{0pt}
\setlength{\parskip}{3pt plus 1pt}
\setlength{\emergencystretch}{3em}   % no overfull text lines
\allowdisplaybreaks                   % long align blocks may break across pages
\clubpenalty=10000 \widowpenalty=10000 \displaywidowpenalty=10000
\raggedbottom                         % no stretched pages before an unbreakable box
\setlist{nosep,leftmargin=1.4em}

\definecolor{accent}{RGB}{22,70,135}
\definecolor{idea}{RGB}{0,120,110}
\definecolor{trap}{RGB}{190,80,0}
\definecolor{ex}{RGB}{60,110,30}
\definecolor{grayline}{RGB}{150,150,150}
\definecolor{c1}{RGB}{22,70,135}
\definecolor{c2}{RGB}{200,90,20}
\definecolor{c3}{RGB}{40,140,60}
\definecolor{c4}{RGB}{170,40,90}

% Source-language term inside output-language text: \term{deadweight loss}
\newcommand{\term}[1]{\textit{\foreignlanguage{\sourcelang}{#1}}}

% --- box titles ---
\newcommand\ttlDefinition{Definition}\newcommand\ttlIntuition{Intuition}\newcommand\ttlTrap{Trap}
\newcommand\ttlMethod{Method}\newcommand\ttlTask{Problem}\newcommand\ttlExample{Example}
\newcommand\ttlKey{Key result}\newcommand\ttlSrcErr{Error in the source}
\newcommand\ttlSolution{Solution.}\newcommand\ttlAnswer{Answer:}
\newcommand\studytitles[9]{%
  \renewcommand\ttlDefinition{#1}\renewcommand\ttlIntuition{#2}\renewcommand\ttlTrap{#3}%
  \renewcommand\ttlMethod{#4}\renewcommand\ttlTask{#5}\renewcommand\ttlExample{#6}%
  \renewcommand\ttlKey{#7}\renewcommand\ttlSrcErr{#8}\renewcommand\ttlSolution{#9}}
\ifdefstring{\studylang}{russian}{\studytitles{Определение}{Интуиция}{Ловушка}{Алгоритм}{Задача}{Пример}{Главное}{Ошибка в исходнике}{Решение.}\renewcommand\ttlAnswer{Ответ:}}{}
\ifdefstring{\studylang}{ngerman}{\studytitles{Definition}{Intuition}{Falle}{Vorgehen}{Aufgabe}{Beispiel}{Merke}{Fehler in der Vorlage}{Lösung.}\renewcommand\ttlAnswer{Ergebnis:}}{}
\ifdefstring{\studylang}{german}{\studytitles{Definition}{Intuition}{Falle}{Vorgehen}{Aufgabe}{Beispiel}{Merke}{Fehler in der Vorlage}{Lösung.}\renewcommand\ttlAnswer{Ergebnis:}}{}
\ifdefstring{\studylang}{french}{\studytitles{Définition}{Intuition}{Piège}{Méthode}{Exercice}{Exemple}{À retenir}{Erreur dans la source}{Solution.}\renewcommand\ttlAnswer{Réponse :}}{}
\ifdefstring{\studylang}{spanish}{\studytitles{Definición}{Intuición}{Trampa}{Método}{Ejercicio}{Ejemplo}{Clave}{Error en la fuente}{Solución.}\renewcommand\ttlAnswer{Respuesta:}}{}

% --- boxes ---
\tcbset{
  base/.style={enhanced,boxrule=0.6pt,arc=2pt,left=6pt,right=6pt,top=4pt,bottom=4pt,
    fonttitle=\bfseries\sffamily\small,before skip=6pt,after skip=6pt},
}
\NewTColorBox{defn}{O{}O{}}{base,colback=accent!4,colframe=accent!70,title={\ttlDefinition#1},#2}
\NewTColorBox{idea}{O{}O{}}{base,colback=idea!5,colframe=idea!75,title={\ttlIntuition#1},#2}
\NewTColorBox{trap}{O{}O{}}{base,colback=trap!5,colframe=trap!80,title={\ttlTrap#1},#2}
\NewTColorBox{method}{O{}O{}}{base,colback=idea!3,colframe=idea!60,title={\ttlMethod#1},#2}
\NewTColorBox{task}{O{}O{}}{base,colback=ex!5,colframe=ex!75,title={\ttlTask#1},#2}
\NewTColorBox{exmpl}{O{}O{}}{base,colback=ex!3,colframe=ex!55,title={\ttlExample#1},#2}
\NewTColorBox{keyres}{O{}O{}}{base,colback=accent!8,colframe=accent,title={\ttlKey#1},#2}
\NewTColorBox{srcerr}{O{}O{}}{base,colback=trap!3,colframe=trap!60,title={\ttlSrcErr#1},#2}
\newcommand{\sol}{\par\smallskip\noindent\textbf{\ttlSolution}\ }
\newcommand{\ans}[1]{\par\smallskip\noindent\colorbox{accent!10}{\parbox{\dimexpr\linewidth-2\fboxsep}{\textbf{\ttlAnswer} #1}}}

% --- headings (never stranded at a page bottom) ---
\titleformat{\section}{\Large\bfseries\sffamily\color{accent}}{\thesection}{0.6em}{}
\titleformat{\subsection}{\large\bfseries\sffamily\color{accent!85!black}}{\thesubsection}{0.5em}{}
\titleformat{\subsubsection}{\normalsize\bfseries\sffamily}{}{0em}{}
\titlespacing*{\section}{0pt}{14pt}{6pt}
\titlespacing*{\subsection}{0pt}{10pt}{4pt}
\titlespacing*{\subsubsection}{0pt}{8pt}{2pt}
% A heading never ends a page: it needs room for itself plus the start of what follows. Before a
% tall unbreakable box, add \needspace{<box height>} yourself or the box moves and leaves the heading.
\pretocmd{\section}{\needspace{8\baselineskip}}{}{}
\pretocmd{\subsection}{\needspace{7\baselineskip}}{}{}
\pretocmd{\subsubsection}{\needspace{5\baselineskip}}{}{}

% --- graphs ---
% \begin{axis}[econ, xlabel=$q$, ylabel=$p$]. Label curves at their end:
%   \addplot[c1,domain=0:8]{2*x+4} node[pos=1,anchor=west,font=\footnotesize]{$MC$};
% With several curves, use pos=0.9/0.95/1 and anchor=south west/west/north west so labels don't collide.
\pgfplotsset{
  every axis/.append style={tick label style={font=\footnotesize},label style={font=\small},
    legend style={font=\footnotesize,draw=none,fill=none}},
  econ/.style={axis lines=left,xlabel style={at={(ticklabel* cs:1)},anchor=north west},
    ylabel style={rotate=-90,at={(ticklabel* cs:1)},anchor=south east},
    xtick=\empty,ytick=\empty,width=6.2cm,height=5cm,samples=120,clip=false,
    every axis plot/.append style={thick},
    axis line style={-{Stealth[length=5pt]}},enlargelimits=false},
}

% Running header: \studyheader{Course · Chapter}{Deliverable}
\newcommand{\studyheader}[2]{%
  \pagestyle{fancy}\fancyhf{}%
  \fancyhead[L]{\small\sffamily\color{grayline}#1}%
  \fancyhead[R]{\small\sffamily\color{grayline}#2}%
  \fancyfoot[C]{\small\thepage}%
  \renewcommand{\headrulewidth}{0.3pt}}
, \renewcommand the \ttl... macros below.
%
% Boxes: defn idea trap method task exmpl keyres srcerr.
%   \begin{trap}[: suffix] ... \end{trap}             title "Trap: suffix"
%   \begin{exmpl}[ 3][breakable] ... \end{exmpl}      long box allowed to split across pages
% Boxes do NOT split across pages by default; pass [breakable] as the 2nd option only for a box
% longer than about half a page.
\providecommand\studylang{english}
\providecommand\sourcelang{english}
\usepackage{etoolbox}
\usepackage[T1,T2A]{fontenc}
\usepackage[utf8]{inputenc}
% Spanish babel's active shorthands (and decimal-dot handling) break TikZ, \% spacing and icomma.
\ifboolexpr{test {\ifdefstring{\studylang}{spanish}} or test {\ifdefstring{\sourcelang}{spanish}}}
  {\PassOptionsToPackage{es-noshorthands,es-nodecimaldot}{babel}}{}
\ifdefstring{\studylang}{\sourcelang}
  {\usepackage[\studylang]{babel}}
  {\usepackage[\sourcelang,\studylang]{babel}}
\usepackage{paratype}
\usepackage{amsmath,amssymb,mathtools}
% Decimal comma for non-English output. In math, type a space after any comma that separates
% items: (1, 2), (x, y), a_{1, 4}. Otherwise (1,2) prints as the decimal 1,2.
\ifdefstring{\studylang}{english}{}{\usepackage{icomma}}
\usepackage{xcolor}
\usepackage[most]{tcolorbox}
\usepackage{enumitem}
\usepackage{booktabs,array,tabularx,multirow}
\usepackage{pgfplots}
\pgfplotsset{compat=1.18}
\usetikzlibrary{arrows.meta,positioning,calc,shapes.geometric,fit,backgrounds,babel}% babel: French/Turkish active : ; ! stay safe in TikZ
\usepackage{titlesec}
\usepackage{fancyhdr}
\usepackage{multicol}
\usepackage{needspace}
\usepackage[hidelinks]{hyperref}

% --- layout defaults that prevent the usual fix-up rounds ---
\setlength{\parindent}{0pt}
\setlength{\parskip}{3pt plus 1pt}
\setlength{\emergencystretch}{3em}   % no overfull text lines
\allowdisplaybreaks                   % long align blocks may break across pages
\clubpenalty=10000 \widowpenalty=10000 \displaywidowpenalty=10000
\raggedbottom                         % no stretched pages before an unbreakable box
\setlist{nosep,leftmargin=1.4em}

\definecolor{accent}{RGB}{22,70,135}
\definecolor{idea}{RGB}{0,120,110}
\definecolor{trap}{RGB}{190,80,0}
\definecolor{ex}{RGB}{60,110,30}
\definecolor{grayline}{RGB}{150,150,150}
\definecolor{c1}{RGB}{22,70,135}
\definecolor{c2}{RGB}{200,90,20}
\definecolor{c3}{RGB}{40,140,60}
\definecolor{c4}{RGB}{170,40,90}

% Source-language term inside output-language text: \term{deadweight loss}
\newcommand{\term}[1]{\textit{\foreignlanguage{\sourcelang}{#1}}}

% --- box titles ---
\newcommand\ttlDefinition{Definition}\newcommand\ttlIntuition{Intuition}\newcommand\ttlTrap{Trap}
\newcommand\ttlMethod{Method}\newcommand\ttlTask{Problem}\newcommand\ttlExample{Example}
\newcommand\ttlKey{Key result}\newcommand\ttlSrcErr{Error in the source}
\newcommand\ttlSolution{Solution.}\newcommand\ttlAnswer{Answer:}
\newcommand\studytitles[9]{%
  \renewcommand\ttlDefinition{#1}\renewcommand\ttlIntuition{#2}\renewcommand\ttlTrap{#3}%
  \renewcommand\ttlMethod{#4}\renewcommand\ttlTask{#5}\renewcommand\ttlExample{#6}%
  \renewcommand\ttlKey{#7}\renewcommand\ttlSrcErr{#8}\renewcommand\ttlSolution{#9}}
\ifdefstring{\studylang}{russian}{\studytitles{Определение}{Интуиция}{Ловушка}{Алгоритм}{Задача}{Пример}{Главное}{Ошибка в исходнике}{Решение.}\renewcommand\ttlAnswer{Ответ:}}{}
\ifdefstring{\studylang}{ngerman}{\studytitles{Definition}{Intuition}{Falle}{Vorgehen}{Aufgabe}{Beispiel}{Merke}{Fehler in der Vorlage}{Lösung.}\renewcommand\ttlAnswer{Ergebnis:}}{}
\ifdefstring{\studylang}{german}{\studytitles{Definition}{Intuition}{Falle}{Vorgehen}{Aufgabe}{Beispiel}{Merke}{Fehler in der Vorlage}{Lösung.}\renewcommand\ttlAnswer{Ergebnis:}}{}
\ifdefstring{\studylang}{french}{\studytitles{Définition}{Intuition}{Piège}{Méthode}{Exercice}{Exemple}{À retenir}{Erreur dans la source}{Solution.}\renewcommand\ttlAnswer{Réponse :}}{}
\ifdefstring{\studylang}{spanish}{\studytitles{Definición}{Intuición}{Trampa}{Método}{Ejercicio}{Ejemplo}{Clave}{Error en la fuente}{Solución.}\renewcommand\ttlAnswer{Respuesta:}}{}

% --- boxes ---
\tcbset{
  base/.style={enhanced,boxrule=0.6pt,arc=2pt,left=6pt,right=6pt,top=4pt,bottom=4pt,
    fonttitle=\bfseries\sffamily\small,before skip=6pt,after skip=6pt},
}
\NewTColorBox{defn}{O{}O{}}{base,colback=accent!4,colframe=accent!70,title={\ttlDefinition#1},#2}
\NewTColorBox{idea}{O{}O{}}{base,colback=idea!5,colframe=idea!75,title={\ttlIntuition#1},#2}
\NewTColorBox{trap}{O{}O{}}{base,colback=trap!5,colframe=trap!80,title={\ttlTrap#1},#2}
\NewTColorBox{method}{O{}O{}}{base,colback=idea!3,colframe=idea!60,title={\ttlMethod#1},#2}
\NewTColorBox{task}{O{}O{}}{base,colback=ex!5,colframe=ex!75,title={\ttlTask#1},#2}
\NewTColorBox{exmpl}{O{}O{}}{base,colback=ex!3,colframe=ex!55,title={\ttlExample#1},#2}
\NewTColorBox{keyres}{O{}O{}}{base,colback=accent!8,colframe=accent,title={\ttlKey#1},#2}
\NewTColorBox{srcerr}{O{}O{}}{base,colback=trap!3,colframe=trap!60,title={\ttlSrcErr#1},#2}
\newcommand{\sol}{\par\smallskip\noindent\textbf{\ttlSolution}\ }
\newcommand{\ans}[1]{\par\smallskip\noindent\colorbox{accent!10}{\parbox{\dimexpr\linewidth-2\fboxsep}{\textbf{\ttlAnswer} #1}}}

% --- headings (never stranded at a page bottom) ---
\titleformat{\section}{\Large\bfseries\sffamily\color{accent}}{\thesection}{0.6em}{}
\titleformat{\subsection}{\large\bfseries\sffamily\color{accent!85!black}}{\thesubsection}{0.5em}{}
\titleformat{\subsubsection}{\normalsize\bfseries\sffamily}{}{0em}{}
\titlespacing*{\section}{0pt}{14pt}{6pt}
\titlespacing*{\subsection}{0pt}{10pt}{4pt}
\titlespacing*{\subsubsection}{0pt}{8pt}{2pt}
% A heading never ends a page: it needs room for itself plus the start of what follows. Before a
% tall unbreakable box, add \needspace{<box height>} yourself or the box moves and leaves the heading.
\pretocmd{\section}{\needspace{8\baselineskip}}{}{}
\pretocmd{\subsection}{\needspace{7\baselineskip}}{}{}
\pretocmd{\subsubsection}{\needspace{5\baselineskip}}{}{}

% --- graphs ---
% \begin{axis}[econ, xlabel=$q$, ylabel=$p$]. Label curves at their end:
%   \addplot[c1,domain=0:8]{2*x+4} node[pos=1,anchor=west,font=\footnotesize]{$MC$};
% With several curves, use pos=0.9/0.95/1 and anchor=south west/west/north west so labels don't collide.
\pgfplotsset{
  every axis/.append style={tick label style={font=\footnotesize},label style={font=\small},
    legend style={font=\footnotesize,draw=none,fill=none}},
  econ/.style={axis lines=left,xlabel style={at={(ticklabel* cs:1)},anchor=north west},
    ylabel style={rotate=-90,at={(ticklabel* cs:1)},anchor=south east},
    xtick=\empty,ytick=\empty,width=6.2cm,height=5cm,samples=120,clip=false,
    every axis plot/.append style={thick},
    axis line style={-{Stealth[length=5pt]}},enlargelimits=false},
}

% Running header: \studyheader{Course · Chapter}{Deliverable}
\newcommand{\studyheader}[2]{%
  \pagestyle{fancy}\fancyhf{}%
  \fancyhead[L]{\small\sffamily\color{grayline}#1}%
  \fancyhead[R]{\small\sffamily\color{grayline}#2}%
  \fancyfoot[C]{\small\thepage}%
  \renewcommand{\headrulewidth}{0.3pt}}
, \renewcommand the \ttl... macros below.
%
% Boxes: defn idea trap method task exmpl keyres srcerr.
%   \begin{trap}[: suffix] ... \end{trap}             title "Trap: suffix"
%   \begin{exmpl}[ 3][breakable] ... \end{exmpl}      long box allowed to split across pages
% Boxes do NOT split across pages by default; pass [breakable] as the 2nd option only for a box
% longer than about half a page.
\providecommand\studylang{english}
\providecommand\sourcelang{english}
\usepackage{etoolbox}
\usepackage[T1,T2A]{fontenc}
\usepackage[utf8]{inputenc}
% Spanish babel's active shorthands (and decimal-dot handling) break TikZ, \% spacing and icomma.
\ifboolexpr{test {\ifdefstring{\studylang}{spanish}} or test {\ifdefstring{\sourcelang}{spanish}}}
  {\PassOptionsToPackage{es-noshorthands,es-nodecimaldot}{babel}}{}
\ifdefstring{\studylang}{\sourcelang}
  {\usepackage[\studylang]{babel}}
  {\usepackage[\sourcelang,\studylang]{babel}}
\usepackage{paratype}
\usepackage{amsmath,amssymb,mathtools}
% Decimal comma for non-English output. In math, type a space after any comma that separates
% items: (1, 2), (x, y), a_{1, 4}. Otherwise (1,2) prints as the decimal 1,2.
\ifdefstring{\studylang}{english}{}{\usepackage{icomma}}
\usepackage{xcolor}
\usepackage[most]{tcolorbox}
\usepackage{enumitem}
\usepackage{booktabs,array,tabularx,multirow}
\usepackage{pgfplots}
\pgfplotsset{compat=1.18}
\usetikzlibrary{arrows.meta,positioning,calc,shapes.geometric,fit,backgrounds,babel}% babel: French/Turkish active : ; ! stay safe in TikZ
\usepackage{titlesec}
\usepackage{fancyhdr}
\usepackage{multicol}
\usepackage{needspace}
\usepackage[hidelinks]{hyperref}

% --- layout defaults that prevent the usual fix-up rounds ---
\setlength{\parindent}{0pt}
\setlength{\parskip}{3pt plus 1pt}
\setlength{\emergencystretch}{3em}   % no overfull text lines
\allowdisplaybreaks                   % long align blocks may break across pages
\clubpenalty=10000 \widowpenalty=10000 \displaywidowpenalty=10000
\raggedbottom                         % no stretched pages before an unbreakable box
\setlist{nosep,leftmargin=1.4em}

\definecolor{accent}{RGB}{22,70,135}
\definecolor{idea}{RGB}{0,120,110}
\definecolor{trap}{RGB}{190,80,0}
\definecolor{ex}{RGB}{60,110,30}
\definecolor{grayline}{RGB}{150,150,150}
\definecolor{c1}{RGB}{22,70,135}
\definecolor{c2}{RGB}{200,90,20}
\definecolor{c3}{RGB}{40,140,60}
\definecolor{c4}{RGB}{170,40,90}

% Source-language term inside output-language text: \term{deadweight loss}
\newcommand{\term}[1]{\textit{\foreignlanguage{\sourcelang}{#1}}}

% --- box titles ---
\newcommand\ttlDefinition{Definition}\newcommand\ttlIntuition{Intuition}\newcommand\ttlTrap{Trap}
\newcommand\ttlMethod{Method}\newcommand\ttlTask{Problem}\newcommand\ttlExample{Example}
\newcommand\ttlKey{Key result}\newcommand\ttlSrcErr{Error in the source}
\newcommand\ttlSolution{Solution.}\newcommand\ttlAnswer{Answer:}
\newcommand\studytitles[9]{%
  \renewcommand\ttlDefinition{#1}\renewcommand\ttlIntuition{#2}\renewcommand\ttlTrap{#3}%
  \renewcommand\ttlMethod{#4}\renewcommand\ttlTask{#5}\renewcommand\ttlExample{#6}%
  \renewcommand\ttlKey{#7}\renewcommand\ttlSrcErr{#8}\renewcommand\ttlSolution{#9}}
\ifdefstring{\studylang}{russian}{\studytitles{Определение}{Интуиция}{Ловушка}{Алгоритм}{Задача}{Пример}{Главное}{Ошибка в исходнике}{Решение.}\renewcommand\ttlAnswer{Ответ:}}{}
\ifdefstring{\studylang}{ngerman}{\studytitles{Definition}{Intuition}{Falle}{Vorgehen}{Aufgabe}{Beispiel}{Merke}{Fehler in der Vorlage}{Lösung.}\renewcommand\ttlAnswer{Ergebnis:}}{}
\ifdefstring{\studylang}{german}{\studytitles{Definition}{Intuition}{Falle}{Vorgehen}{Aufgabe}{Beispiel}{Merke}{Fehler in der Vorlage}{Lösung.}\renewcommand\ttlAnswer{Ergebnis:}}{}
\ifdefstring{\studylang}{french}{\studytitles{Définition}{Intuition}{Piège}{Méthode}{Exercice}{Exemple}{À retenir}{Erreur dans la source}{Solution.}\renewcommand\ttlAnswer{Réponse :}}{}
\ifdefstring{\studylang}{spanish}{\studytitles{Definición}{Intuición}{Trampa}{Método}{Ejercicio}{Ejemplo}{Clave}{Error en la fuente}{Solución.}\renewcommand\ttlAnswer{Respuesta:}}{}

% --- boxes ---
\tcbset{
  base/.style={enhanced,boxrule=0.6pt,arc=2pt,left=6pt,right=6pt,top=4pt,bottom=4pt,
    fonttitle=\bfseries\sffamily\small,before skip=6pt,after skip=6pt},
}
\NewTColorBox{defn}{O{}O{}}{base,colback=accent!4,colframe=accent!70,title={\ttlDefinition#1},#2}
\NewTColorBox{idea}{O{}O{}}{base,colback=idea!5,colframe=idea!75,title={\ttlIntuition#1},#2}
\NewTColorBox{trap}{O{}O{}}{base,colback=trap!5,colframe=trap!80,title={\ttlTrap#1},#2}
\NewTColorBox{method}{O{}O{}}{base,colback=idea!3,colframe=idea!60,title={\ttlMethod#1},#2}
\NewTColorBox{task}{O{}O{}}{base,colback=ex!5,colframe=ex!75,title={\ttlTask#1},#2}
\NewTColorBox{exmpl}{O{}O{}}{base,colback=ex!3,colframe=ex!55,title={\ttlExample#1},#2}
\NewTColorBox{keyres}{O{}O{}}{base,colback=accent!8,colframe=accent,title={\ttlKey#1},#2}
\NewTColorBox{srcerr}{O{}O{}}{base,colback=trap!3,colframe=trap!60,title={\ttlSrcErr#1},#2}
\newcommand{\sol}{\par\smallskip\noindent\textbf{\ttlSolution}\ }
\newcommand{\ans}[1]{\par\smallskip\noindent\colorbox{accent!10}{\parbox{\dimexpr\linewidth-2\fboxsep}{\textbf{\ttlAnswer} #1}}}

% --- headings (never stranded at a page bottom) ---
\titleformat{\section}{\Large\bfseries\sffamily\color{accent}}{\thesection}{0.6em}{}
\titleformat{\subsection}{\large\bfseries\sffamily\color{accent!85!black}}{\thesubsection}{0.5em}{}
\titleformat{\subsubsection}{\normalsize\bfseries\sffamily}{}{0em}{}
\titlespacing*{\section}{0pt}{14pt}{6pt}
\titlespacing*{\subsection}{0pt}{10pt}{4pt}
\titlespacing*{\subsubsection}{0pt}{8pt}{2pt}
% A heading never ends a page: it needs room for itself plus the start of what follows. Before a
% tall unbreakable box, add \needspace{<box height>} yourself or the box moves and leaves the heading.
\pretocmd{\section}{\needspace{8\baselineskip}}{}{}
\pretocmd{\subsection}{\needspace{7\baselineskip}}{}{}
\pretocmd{\subsubsection}{\needspace{5\baselineskip}}{}{}

% --- graphs ---
% \begin{axis}[econ, xlabel=$q$, ylabel=$p$]. Label curves at their end:
%   \addplot[c1,domain=0:8]{2*x+4} node[pos=1,anchor=west,font=\footnotesize]{$MC$};
% With several curves, use pos=0.9/0.95/1 and anchor=south west/west/north west so labels don't collide.
\pgfplotsset{
  every axis/.append style={tick label style={font=\footnotesize},label style={font=\small},
    legend style={font=\footnotesize,draw=none,fill=none}},
  econ/.style={axis lines=left,xlabel style={at={(ticklabel* cs:1)},anchor=north west},
    ylabel style={rotate=-90,at={(ticklabel* cs:1)},anchor=south east},
    xtick=\empty,ytick=\empty,width=6.2cm,height=5cm,samples=120,clip=false,
    every axis plot/.append style={thick},
    axis line style={-{Stealth[length=5pt]}},enlargelimits=false},
}

% Running header: \studyheader{Course · Chapter}{Deliverable}
\newcommand{\studyheader}[2]{%
  \pagestyle{fancy}\fancyhf{}%
  \fancyhead[L]{\small\sffamily\color{grayline}#1}%
  \fancyhead[R]{\small\sffamily\color{grayline}#2}%
  \fancyfoot[C]{\small\thepage}%
  \renewcommand{\headrulewidth}{0.3pt}}
